\documentclass[11pt]{article}

\usepackage[margin=1in]{geometry}
\usepackage{amsmath,amssymb,amsthm}
\usepackage{mathtools}
\usepackage{enumitem}
\usepackage[hidelinks]{hyperref}
\usepackage[expansion=false]{microtype}
\usepackage{xr}


\newtheorem{theorem}{Theorem}[section]
\newtheorem{lemma}[theorem]{Lemma}
\newtheorem{proposition}[theorem]{Proposition}
\newtheorem{corollary}[theorem]{Corollary}
\newtheorem{conjecture}[theorem]{Conjecture}
\theoremstyle{definition}
\newtheorem{definition}[theorem]{Definition}
\newtheorem{hypothesis}[theorem]{Hypothesis}
\newtheorem{remark}[theorem]{Remark}

\DeclareMathOperator{\lcm}{lcm}
\newcommand{\NN}{\mathbb{N}}
\newcommand{\ZZ}{\mathbb{Z}}
\newcommand{\QQ}{\mathbb{Q}}
\newcommand{\PP}{\mathcal{P}}
\newcommand{\der}{{}'}
\newcommand{\dd}{\,\mathrm{d}}

\title{Obstructions for Erd\H{o}s Problem~\#307:\\ places, potentials, and two closures on the existence route}
\author{Vico Bonfioli\thanks{\texttt{vicobonfioli@gmail.com}.
  See the acknowledgments for a candid note on the use of AI assistance in this work.}}
\date{Version 1.75.0\quad 2 September 2026}

\externaldocument[K-]{erdos307-core}
\externaldocument[A-]{erdos307-computational}
\externaldocument[C-]{erdos307-analytic}
\begin{document}
\maketitle

\begin{abstract}
This note records what cannot be used on Erd\H{o}s Problem~\#307, in the form of theorems rather
than of failed attempts. Over $\mathbb{F}_q[t]$ the derivative admits no two--cycle, by a degree
argument with no integer counterpart, and no place of $\mathbb{Q}$ carries that argument: the
archimedean absolute value is expanding on the relevant range, and each $p$--adic valuation, while
lowered at the primes of $n$, is raised elsewhere. By Ostrowski these are all the places, so the
function--field mechanism is unavailable as a classification statement and not as a failure of
ingenuity. Twisted two--cycles do exist over $\mathbb{Z}[i]$, $\mathbb{Z}[\sqrt{-2}]$,
$\mathbb{Z}[\omega]$ and $\mathbb{Z}[\sqrt2]$, so the obstruction is the orderability of
$\mathbb{Q}$. Two closures on the existence route are proved: no certificate distinguishes the two
branches, and no potential of the form $\log n+G(\sigma(n))$ can decrease around a cycle, for any
$G$, since the $G$--terms telescope and only the mass identity survives.
\end{abstract}
\section{A function--field lens: the obstruction is archimedean}\label{sec:ff}

Define an arithmetic derivative on $\mathbb{F}_q[t]$ by $P'=1$ for monic irreducibles and the Leibniz
rule, so $f'=\sum_i e_i\,f/P_i$ for $f=\prod_i P_i^{e_i}$ (the exponents $e_i$ read modulo the
characteristic).

\begin{theorem}\label{thm:ff}
Over $\mathbb{F}_q[t]$ this derivative has no two--cycles and no nonzero fixed points. Indeed
$\deg f'\le\deg f-1$ for every nonconstant $f$. The analogues of the integer fixed points $p^{p}$ also
vanish: $(P^{p})'=p\,P^{p-1}P'=0$ in characteristic $p$. The same holds verbatim for every twisted
variant ($P'=u_P\in\mathbb{F}_q^\times$, Leibniz).
\end{theorem}
\begin{proof}
Each term $e_i\,f/P_i$ has degree $\deg f-\deg P_i\le\deg f-1$, so $\deg f'\le\deg f-1$. A two--cycle
$f'=g$, $g'=f$ would force $\deg g\le\deg f-1$ and $\deg f\le\deg g-1$, hence $\deg f\le\deg f-2$;
a fixed point $f'=f$ is excluded likewise. (Verified exhaustively over $\mathbb{F}_2$ to degree $7$,
$\mathbb{F}_3$ to degree $5$, and $\mathbb{F}_5$ to degree $4$; twisted variants over $\mathbb{F}_2$ and
$\mathbb{F}_5$ on small supports through $2.3\times10^{6}$ twist configurations, none closing.)
\end{proof}

\noindent\textup{Formalised.} \texttt{Erdos307.ff\_no\_cycle} is the theorem, on the single
hypothesis $\deg f'+1\le\deg f$; \texttt{deg\_drop\_of\_terms} is that hypothesis from the term
bound, and \texttt{no\_fixed\_point\_of\_drop} excludes fixed points. The argument is stated on the
degrees alone, so it holds verbatim for every twisted variant, the twist multiplying each term by a
unit \textup{(\texttt{lean/Erdos307/FunctionField.lean})}.

\begin{corollary}[The function--field Erd\H{o}s--Barbeau equation is insoluble]\label{cor:ff307}
For every finite field $\mathbb{F}_q$, all nonempty finite sets $P,Q$ of monic irreducibles, and all
unit weights $u_\pi,v_\rho\in\mathbb{F}_q^\times$,
\[
\Bigl(\sum_{\pi\in P}\frac{u_\pi}{\pi}\Bigr)\Bigl(\sum_{\rho\in Q}\frac{v_\rho}{\rho}\Bigr)\ne1.
\]
\end{corollary}
\begin{proof}
Write $f=\prod P$, $g=\prod Q$, $D_u(f)=\sum_\pi u_\pi f/\pi$; the equation reads $D_u(f)\,D_v(g)=fg$.
If some $\pi\in P\cap Q$, then $\pi^2\mid fg$ while $D_u(f)\equiv u_\pi f/\pi\not\equiv0$ and
$D_v(g)\equiv v_\pi g/\pi\not\equiv0\pmod\pi$, so $\pi\nmid D_u(f)D_v(g)$; impossible; hence
$\gcd(f,g)=1$. The same congruence gives $\gcd(f,D_u(f))=\gcd(g,D_v(g))=1$, so $g\mid D_u(f)$ and
$f\mid D_v(g)$, forcing $D_u(f)=c\,g$, $D_v(g)=c^{-1}f$ with $c\in\mathbb{F}_q^\times$: a twisted
two--cycle, excluded by the degree drop. The weighted product equation is thus insoluble over every
$\mathbb{F}_q(t)$; it \emph{is} soluble over the unit--rich number rings of
Theorem~\ref{thm:existence}, since any twisted cycle gives $(D_u(a)/a)(D_v(b)/b)=(b/a)(a/b)=1$, and over
$\mathbb{Q}$, where the admissible weights are $\pm1$, it is \#307 itself with the signed barrier of
Remark~\ref{rem:ordered}, open beyond $2.09\times10^{56}$. The three regimes (non--archimedean
monotonicity (closed, negative), unit--broken monotonicity (closed, positive), ordered--archimedean
(open, barred)) now have theorems on two of three doors.
\end{proof}

Over $\mathbb{Z}$, by contrast, $n'=n\sum_i e_i/p_i$ can exceed $n$ precisely because the archimedean
absolute value lets a sum of strictly smaller terms outgrow the original; the ``mass $>1$''
phenomenon $\sum 1/p>1$, which has no degree--theoretic analogue. Theorem~\ref{thm:ff} thus localises
the entire existence question to the archimedean place: every non--archimedean shadow of \#307 is
trivially negative. This is consistent with (and partly explains) both the census finding of
Section~\ref{K-sec:computations} of the core note and Proposition~\ref{A-prop:nolocal} (no single modulus obstructs the
cross--match system in the tested range); the difficulty is not local. Any eventual proof must therefore distinguish $\mathbb{Z}$ from $\mathbb{F}_q[t]$; over the
latter the decisive tools are the degree and the Mason--Stothers/Wronskian inequality
\cite{Stothers}, neither of which has a $\mathbb{Z}$--analogue. The conjectural analogue, $abc$,
does not fill the gap, and quantifiably so. Every cycle inherits, from $N'=a^2+b^2$ with $N=ab$, the
relation $(a-b)^2+4ab=(a+b)^2$ on coprime squarefree members, but this triple is $abc$--unexceptional,
of quality $q=\log(a+b)^2/\log\operatorname{rad}\!\bigl((a-b)\cdot ab\cdot(a+b)\bigr)\approx\tfrac12$
(median $0.543$, never above $0.74$ across $6\times10^4$ coprime squarefree pairs; the only $q>1$ cases
are degenerate smooth--gap pairs such as $(5,3)$, vanishingly far below the $10^{56}$ scale a solution
must occupy). Granting $abc$ in its strong explicit form would therefore yield no bound here: the
function--field collapse rests on an inequality genuinely \emph{absent} over $\mathbb{Z}$, not merely
unproven, and the size mismatch is by a factor of two in the exponent. \#307 lives in the
$abc$--comfortable zone.

\emph{The free grading, and what evaluation destroys.} Theorem~\ref{thm:ff} has a structural reading. In the free Leibniz world (the polynomial ring
$\mathbb{Z}[x_p]$ with $x_p'=1$) the derivative of a squarefree monomial of degree $n$ is the
elementary symmetric form $e_{n-1}$, homogeneous of degree $n-1$: the formal derivative strictly lowers
the grading, cycles are impossible, and over $\mathbb{F}_q[t]$ the degree realises this grading
faithfully. Over $\mathbb{Z}$ the evaluation $x_p\mapsto p$ destroys homogeneity, and the data show it
destroys it completely: for random squarefree $a$ of fixed size, $\omega(a')$ is governed by the size of
$a'$ alone (Erd\H{o}s--Kac), flat in $\omega(a)$ within each parity class (measured: mean
$\omega(a')=2.7$--$3.0$ as $\omega(a)$ runs over $3,5,7$, and $3.6$--$4.0$ over $4,6,8$, the offset being
the parity anatomy of Theorem~\ref{K-thm:parity}). A two--cycle is therefore a doubly anti--typical event:
an additively assembled integer must carry multiplicative rank far in the Erd\H{o}s--Kac tail \emph{and}
on an exactly prescribed support; the two independent prices whose product is the heuristic rarity of
Section~\ref{A-sec:anatomy} of the computational companion. The deeper additive--multiplicative content collapses, on inspection, into
the $S$--unit form of Proposition~\ref{prop:sunit}: dividing the cycle equation by $D_Q$ exhibits it as
an $n$--term unit equation; one more coordinate system in which the same exactness reappears
unchanged.


\begin{remark}[The obstruction is \emph{ordered}--archimedean]\label{rem:ordered}
The barrier is finer than ``archimedean'': its engine is AM--GM on real, \emph{positive} masses, i.e.\
the ordering of $\mathbb{R}$. Over $\mathbb{Z}[i]$ (a PID, so the Rigidity Lemma and cross--matching
survive verbatim, with cycles defined up to units), the masses $\sum1/\pi$ are \emph{complex}: phases can
align so that $\bigl|\sum1/\pi\bigr|>1$ with as few as three Gaussian primes (e.g.\
$\tfrac1{1+i}+\tfrac1{2+i}+\tfrac1{2-i}=1.3-0.5i$), and the $59$--prime bound has no analogue. The
ordered--archimedean thesis thus predicts that small derivative two--cycles \emph{may} exist over
$\mathbb{Z}[i]$ where they are barred over $\mathbb{Z}$, and the prediction is confirmed. Two facts
make the contrast precise.
\emph{First}, phases do not help over $\mathbb{Z}$: for any \emph{signed} derivative ($p'=\pm1$, Leibniz)
a two--cycle still satisfies $|s||t|=1$ with $|s|\le\sigma(a)$, $|t|\le\sigma(b)$ and disjoint supports
(the rigidity proof is unchanged), so $\sigma(a)+\sigma(b)\ge2$ and the full $59$--prime,
$2.09\times10^{56}$ barrier applies to every choice of signs: real phases cannot cancel.
\emph{Second}, fourth roots of unity suffice instantly: Leibniz derivatives on $\mathbb{Z}[i]$ with unit
prime--values ($\pi'\in\{\pm1,\pm i\}$) possess genuine two--cycles at norm ${\sim}10^4$. The smallest:
\[ a=(1+i)(2+7i)(6+11i)=-129-i,\qquad b=3\,(2+i)(1+2i)(6+5i)=-75+90i, \]
with $a'=b$ under the \emph{standard} values $\pi'=1$ on the primes of $a$, and $b'=a$ exactly under
$\pi'=i$ on $3$, $2+i$, $1+2i$ and $\pi'=1$ on $6+5i$; seven prime classes
($N(a)=16{,}642$, $N(b)=13{,}725$) against the fifty--nine that $\mathbb{Z}$ demands. Exhaustive search
finds exactly sixteen such cycles up to conjugation and direction having a member of norm
$\le2\times10^7$ with $\omega\le5$ on the enumerated side ($\omega$--profiles $(3,4)$, $(4,4)$,
$(5,4)$, and even $(4,6)$) and the population grows steadily ($4$ below $3\times10^5$, $16$ below
$2\times10^7$), consistent with an infinite, slowly growing family; in six of the sixteen one side
requires no twist at all. All sixteen were verified exactly, by independent recomputation; no twisted
\emph{fixed} point ($a'=ua$) exists in the same range. The strict two--sided normalisation $\pi'=1$
(first--quadrant representatives) has no cycle up to units (nor any fixed point, zero--derivative, or
unit--derivative element) with $N(a)\le2\times10^9$ ($9.4\times10^8$ candidates, partners by full
factorisation); this is the expected order rather than a separate rigidity, a single fixed phase system
being heuristically $4^{\omega-1}$--fold rarer than the union of all systems. The contrast with
$\mathbb{Z}$ survives this accounting: over $\mathbb{Z}$ \emph{every} one of the $2^{\omega-1}$ sign
systems is barred by the theorem above, while over $\mathbb{Z}[i]$ the union of phase systems is
demonstrably populated from norm $1.7\times10^4$ on. The barrier of \#307 is thus localised precisely: it is the impossibility of phase
cancellation in the ordered field $\mathbb{R}$, where the only available phases $\pm1$ provably gain
nothing; while the moment fourth roots of unity are admitted, the analogous existence problem closes
at seven primes. The cycles found come in conjugate pairs (supports not conjugation--stable); a
conjugation--\emph{stable} Gaussian cycle would descend to rational data through the norm form
$a^2+b^2$, i.e.\ precisely the Pythagorean layer of Proposition~\ref{K-prop:pairform}; now a concrete
target rather than a speculation.
\end{remark}

\begin{remark}[The descended operator; the symmetric sector]\label{rem:descent}
Conjugation--stable supports admit an exact descent to $\mathbb{Z}$. Such a squarefree Gaussian integer
is an associate of an odd squarefree rational $m$ (split primes entering as whole conjugate pairs; the
$(1+i)$--sector is empty by a short parity argument), and a conjugation--equivariant unit assignment
makes the twisted derivative rational:
\[ D(m)\;=\;\sum_{\substack{p\mid m\\ p\equiv1\,(4)}} t_p\,\frac mp\;+\;\sum_{\substack{q\mid m\\
q\equiv3\,(4)}}\varepsilon_q\,\frac mq, \qquad t_p\in\{\pm2x_p,\pm2y_p\},\quad \varepsilon_q=\pm1, \]
where $p=x_p^2+y_p^2$; a Leibniz derivative on odd squarefree integers whose prime values have size
$2\sqrt p$ rather than $1$, and whose masses ${\sim}2/\sqrt p$ reach $2$ with four primes, so no
$59$--prime barrier exists for it. A two--cycle of $D$ is exactly a conjugation--stable Gaussian cycle
and would plant the phenomenon on rational soil. Parity constrains the hunt: $D(m)$ is odd iff $m$
carries an odd number of primes $\equiv3\pmod4$, so every member must contain such a prime.
One level deeper, since $q^{-1}\equiv q\pmod 8$ for odd $q$, the descended operator obeys
$D(m)\equiv m\,W(m)\pmod 8$ with $W(m)=\sum_{p\mid m}t_p\,p$, so a two--cycle forces
$W(m)\equiv W(m')\equiv mm'\pmod 8$; a mod--$8$ constraint on the sign assignment that prunes
the $4^{\#\mathrm{split}}2^{\#\mathrm{inert}}$ phase enumeration severalfold in future searches
(verified on $2\times10^5$ random assignments; \texttt{code/peeling\_groupoid.py}). An
exhaustive search over members $m\le10^9$ (Gaussian norm $10^{18}$, $\omega\le8$, $1.4\times10^9$
closure tests) found no cycle, no fixed point $|D(m)|=m$, and no vanishing $D(m)=0$. The search
has since been extended two decades further in $m$, on a box narrower in the other parameters:
$\omega\in[4,6]$ and odd primes below $600$, with $m$ running over $[10^9,10^{10}]$ and then
$[10^{10},10^{11}]$. Both tiers are exhaustive over their $5{,}778$ prime pairs and return nothing:
$10{,}004{,}024$ members and $6.94\times10^8$ closure tests on the first, $22{,}218{,}430$ members
and $1.82\times10^9$ closure tests on the second, so $3.22\times10^7$ members and
$2.52\times10^9$ closure tests in all, with no two--cycle at any point
\textup{(\texttt{code/hunt/descended\_hunt2.rs}, \texttt{code/hunt/run\_tier2\_par.sh})}. The tier boxes are
not nested in the one above, being narrower in $\omega$ and in the prime range while wider in $m$,
so the two counts complement rather than supersede one another; and the tiers test closure only, not
the fixed--point and vanishing conditions recorded above. We claim no
obstruction from this: equivariance prunes the available phases fourfold per split prime, and a phase
count makes emptiness at this depth the expected order. The descent target remains open over
$\mathbb{Z}[i]$, but not in the real--quadratic world, where the corresponding sector contains a fully
rational pair (Remark~\ref{rem:orderable}). Finally, the
descended coefficients are classical objects: at a split prime $p$ the four values
$\{\pm2x_p,\pm2y_p\}$ are exactly the Frobenius traces of the congruent--number curve $y^2=x^3-x$
(CM by $\mathbb{Z}[i]$) and of its twists ($a_p\in\{\pm2x_p,\pm2y_p\}$ verified by point counting
for all $p<120$) so the conjugation--stable sector is the arithmetic derivative twisted by the
Hecke eigensystem of $\mathbb{Q}(i)$, and its mass statistics are governed unconditionally by Hecke's
equidistribution theorem: the first point of contact between \#307 and automorphic forms. Concretely:
over random squarefree $m$ the three--series conditions converge (the variance term is $\sum 2/p^2$), so
the descended mass $D(m)/m$ has an Erd\H{o}s--Wintner \emph{limiting distribution} whose prime components
are arcsine laws scaled by $1/\sqrt p$; the required input, equidistribution of the Hecke angles, is
visible exactly (Kolmogorov--Smirnov distance $0.0041$ from uniform over the $8{,}977$ split primes below
$2\times10^5$, well inside the $95\%$ band $0.0144$).
\end{remark}

\begin{remark}[The wall is the orderability of $\mathbb{Q}$]\label{rem:orderable}
The cycles over $\mathbb{Z}[i]$ permit a dissection of the barrier that no amount of theory over
$\mathbb{Z}$ could provide. The entire structural chain of this note (rigidity, cross--matching, the
Pythagorean identity, the split discriminant) transfers verbatim to twisted derivatives over
$\mathbb{Z}[i]$: Leibniz holds for disjoint supports, so a cycle $(a,b)$ has $D(N)=a^2+b^2$ for $N=ab$,
with $D(N)-2N=(a-b)^2$ and $2N-D(N)=\bigl(i(a-b)\bigr)^2$ (all verified exactly on the cycles found).
The last equation exposes the wall: over $\mathbb{Z}[i]$, where $-1=i^2$ is a square, \emph{both}
discriminant factors are squares automatically, and the plus/minus distinction of
Proposition~\ref{C-prop:splitdisc} (which over $\mathbb{Z}$ carries the entire obstruction, via
$(a-b)^2\ge0\Rightarrow\sigma(N)\ge2$) has no content. The unique non--transferring link is the
sign--definiteness of squares: the barrier is, in the end, the \emph{formal reality} of $\mathbb{Q}$ in
the sense of Artin--Schreier. (The two Artin--Schreier theorems thus bracket the problem: their
characteristic--$2$ map builds the tower of Proposition~\ref{C-prop:ff-pyth}, their theory of orderable
fields names the wall.) This classifies where the mechanism can operate at all: it needs an ordering
making squares nonnegative (formal reality, hence a real embedding) \emph{and} bounded twisting phases
(finite unit group); by Dirichlet's unit theorem, unit rank $r_1+r_2-1=0$ with $r_1\ge1$ forces degree
one. \emph{$\mathbb{Q}$ is the only number field over which the mass barrier can operate}: over
imaginary quadratic fields squares lose sign--definiteness (over $\mathbb{Z}[i]$ even $-1$ is a square),
and over every other field the unit rank is positive and twisted masses are unbounded. This is a
classification of the wall's habitat, not an existence claim elsewhere, but over $\mathbb{Z}[i]$
existence is no longer a claim: it is a list of sixteen. The ladder of imaginary quadratic rings
confirms the mechanism quantitatively. Over $\mathbb{Z}[\sqrt{-2}]$, whose only units are $\pm1$,
precisely the twisting freedom that provably gains nothing over $\mathbb{Z}$; sign--twisted cycles
nevertheless exist from norm $11{,}462$ on ($a=-90-41\sqrt{-2}$, $b=75-45\sqrt{-2}$, seven prime
classes): the complex embedding alone carries the phenomenon. Over the Eisenstein ring
$\mathbb{Z}[\omega]$, with six units, cycles appear at norm $2{,}964$ with \emph{six} prime classes
($a=-40+22\omega$, $b=56-21\omega$); the smallest derivative two--cycle exhibited in any ring,
and the population at fixed height grows with the unit group: at norm $\le2\times10^7$ the censuses
read $14$, $56$, and $278$ twisted--cycle instances for unit groups of order $2$, $4$, $6$ (a factor
${\sim}4$--$5$ per two extra units, as a phase count predicts) while the strict canonical systems of
all three rings remain empty to norm $2\times10^9$ ($3\times10^9$ candidates in total, partners by full
factorisation). Two further objects appeared over $\mathbb{Z}[\sqrt{-2}]$: a composite with \emph{unit}
derivative ($a'\sim1$; impossible over $\mathbb{Z}$, where $n'=1$ characterises the primes), and an
\emph{anti--holomorphic} cycle: $a=-3163-1474\sqrt{-2}$ satisfies $a'=-\bar a$ exactly, under a single
conjugation--equivariant sign assignment closing both directions ($N(a)=N(b)=14{,}349{,}921$, five prime
classes each side). The equation $a'=u\bar a$ is vacuous over $\mathbb{Z}$ only in its
\emph{one--dimensional} reading, where it says $n'=\pm n$. In the reading that matches the problem it is
not vacuous but central: by Proposition~\ref{C-prop:multiplier} the Pythagorean layer over $\mathbb{Z}$
\emph{is} the anti--holomorphic equation $\mathcal{D}z=\mu\bar z$ for the coordinatewise derivative on
$z=a+bi$, and \#307 is exactly its unit--multiplier case. The ring ladder and the main problem are
therefore instances of one equation rather than analogues; what singles out $\mathbb{Z}$ is the extra
constraint $\operatorname{Im}\mu=1$, which confines the multiplier to a line meeting
$\mathbb{Z}[i]^{\times}$ in the single point $\mu=i$. That the anti--holomorphic sector is thin is
uniform across the ladder: over $\mathbb{Z}[\sqrt{-2}]$ the cycle above is still the \emph{only} one
after widening the census simultaneously in all three directions; prime--norm $\le2\times10^4$,
$\omega\le8$, product--norm $\le10^{10}$ (a factor $700$ past the cycle itself), $8.09\times10^8$
supports, \texttt{code/hunt/antiholo\_hunt.rs}, and over
$\mathbb{Z}$ the sector is empty below $10^{112}$ by Corollary~\ref{K-cor:emptytest}. We record the
plateau as consistent with genuine finiteness, not as evidence for it: a random model for
$a'=-\bar a$ predicts a count growing like $\log X$, which the tested range cannot separate from a
constant. All cycles were verified exactly by independent recomputation. Finally, the second pillar was tested in isolation: over
$\mathbb{Z}[\sqrt2]$ (\emph{totally real}, so squares are sign--definite in both embeddings, but with
infinite unit group $\pm(1+\sqrt2)^{\mathbb{Z}}$) twisted cycles appear already at norm $882$: with
unit twists of height $\le2$, $a=21\sqrt2$ (classes $\sqrt2$, $-1+2\sqrt2$, $1+2\sqrt2$, $3$) cycles
exactly with $b=51+13\sqrt2$, of norm $2263$ and only \emph{two} prime classes (against the
fifty--nine that $\mathbb{Z}$ demands) and the composite $369-15\sqrt2$ has derivative exactly
$\varepsilon^6=99+70\sqrt2$, a unit. Both pillars of the classification are thus load--bearing, and each
is now confirmed by explicit counterexample on the soil where it alone fails: drop formal reality and
cycles appear (imaginary quadratic rings, norms $10^3$--$10^4$); keep formal reality but drop unit
finiteness and cycles appear (real quadratic, norm $882$, the smallest in any ring); over $\mathbb{Q}$,
where both hold, the barrier stands at $10^{112}$ by theorem. The bracketing of $\mathbb{Q}$ is complete
from both axes. The census sharpens it: at norm $\le2\times10^7$ the real--quadratic ring carries
$30{,}204$ twisted--cycle instances against $14$, $56$, $278$ for the imaginary rings with $2,4,6$ units
infinite unit rank dominates by two orders of magnitude, and contains both the \emph{minimal}
architecture, $a=\sqrt2\,(5-\sqrt2)$ cycling with $b=7$ at norms $(46,49)$, two prime classes per member
(the rational prime $7$ being composite in the ring), the smallest derivative cycle found in any ring;
and a \emph{fully rational} pair:
\[ D(3707)=1547,\qquad D(1547)=3707\qquad(3707=11\cdot337,\ \ 1547=7\cdot13\cdot17), \]
exactly, with unit heights $\le2$ on both sides, verified independently: two ordinary integers in a
derivative two--cycle. The descent phenomenon that is empty over $\mathbb{Z}[i]$ to norm $10^{18}$ is
realised over $\mathbb{Z}[\sqrt2]$ at four digits. The canonical system of $\mathbb{Z}[\sqrt2]$ remains
cycle--free to norm $2\times10^9$ ($7.7\times10^8$ candidates), now with seven unit--derivative
composites.
\end{remark}

The cycle phenomenon over rings with units is governed by the theory of unit equations, which supplies
its rigidity theorem.

\begin{proposition}[$S$--unit rigidity of twisted cycles]\label{prop:sunit}
Let $\mathcal{O}$ be the ring of integers of a number field and $S$ a finite set of primes of
$\mathcal{O}$. If $(a,b)$ is a twisted two--cycle with $\operatorname{supp}(ab)\subseteq S$, then
dividing $D(a)=b$ by $b$ exhibits a solution of
\[ x_1+\dots+x_{\omega(a)}=1,\qquad x_i=u_i\,\frac{a/\pi_i}{b}, \]
in the $S$--unit group of the field, and likewise for $b$. By the theorem of
Evertse--Schlickewei--Schmidt \cite{ESS2002} (cf.\ \cite{Evertse1984}), such an equation has only
finitely many nondegenerate solutions; hence \emph{every fixed support carries at most finitely many
twisted cycles free of vanishing subsums}, and the species can be infinite only through growing
supports; precisely the behaviour of every census above. (Vanishing subsums correspond to partial
zero--derivative patterns; none occurred in any search. The relation was verified exactly on the
$\mathbb{Z}[\sqrt2]$ specimen: the four $S$--unit terms of $a=21\sqrt2$ sum to $1$.)
\end{proposition}

\begin{remark}[Galois folding, and the thinness of symmetric sectors]\label{rem:folding}
Over a quadratic field with automorphism $\sigma$, a $\sigma$--equivariant choice of unit values
($u_{\pi^\sigma}=u_\pi^{\,\sigma}$) makes $D$ commute with conjugation, so the \emph{single} equation
$D(a)=a^\sigma$ already implies that $(a,a^\sigma)$ is a twisted two--cycle (checked symbolically and on
$1{,}988$ random instances). In the two--prime case the partner is forced:
$\rho=-\bigl(\bar u v\,\pi+\bar v\,(\pi^\sigma)^2\bigr)/\bigl(N(u)-N(\pi)\bigr)$, with integrality a
congruence and the entire residue of the problem the primality of $|N(\rho)|$ along Pell--exponentially
growing sequences; Lehmer territory rather than Bunyakovsky. Whether a folded sector is starved or
rich is decided by its degrees of freedom. When the fold leaves no slack the sector is thin: the folded
\emph{two}--prime family over $\mathbb{Z}[\sqrt2]$ is empty in the tested box ($36{,}000$ parameter
triples, none even integral), the conjugation--stable sector over $\mathbb{Z}[i]$ is empty to Gaussian
norm $10^{18}$, and the anti--holomorphic cycle over $\mathbb{Z}[\sqrt{-2}]$ is unique in its census.
But with \emph{three} primes the folded equation
$u_1\pi_2\pi_3+u_2\pi_1\pi_3+u_3\pi_1\pi_2=(\pi_1\pi_2\pi_3)^\sigma$ carries three unit unknowns against
two real equations, and this underdetermined system is \emph{populated}: an exhaustive box ($p_i\le137$,
unit heights $\le4$) contains twenty--four solutions over eleven prime triples, the smallest being
\[ a=(3+\sqrt2)(5-2\sqrt2)(5-\sqrt2)=57-16\sqrt2,\qquad D(a)=a^\sigma=57+16\sqrt2, \]
with unit values $\varepsilon^2$, $\varepsilon$, $-\varepsilon^{-1}$: \emph{the derivative conjugates the
element}, at norm $2737$ on both sides, the equivariant return being automatic (verified exactly).
Existence of twisted cycles over the four rings tested is a theorem by exhibition; for \emph{every} ring
with infinite unit group it is now reduced to the abundance of lattice points on these underdetermined
unit systems (measured plentiful already at $k=3$) a geometry--of--numbers question rather than a
primality one.
\end{remark}

\begin{remark}[{The forced third prime: a Schinzel--type reduction over $\mathbb{Z}[\sqrt2]$, unblocked}]
\label{rem:forced}
Within that abundance the arithmetic is sharper than free lattice--counting. Fixing $\pi_1,\pi_2$ and the
units $u_1,u_2,u_3$, the folded equation rearranges to
$\pi_3\,(u_1\pi_2+u_2\pi_1)-(\pi_1\pi_2)^\sigma\pi_3^\sigma=-u_3\pi_1\pi_2$, a $2\times2$ integer linear
system for $\pi_3=(x,y)$ whose determinant is $N(u_1\pi_2+u_2\pi_1)-N(\pi_1\pi_2)$ (an identity, verified
on all $168$ solutions of the box). So $\pi_3$ is \emph{forced} by $(\pi_1,\pi_2,u_1,u_2,u_3)$, and
infinitude of twisted cycles is exactly the primality of this forced $\pi_3$ for infinitely many such data
a Schinzel/Bunyakovsky--type condition over $\mathbb{Z}[\sqrt2]$. Crucially it is \emph{not} the regime
of Proposition~\ref{K-prop:nopoly}: the exponential unit freedom $u_i=\pm\varepsilon^{k_i}$ keeps the family
non--polynomial yet populated, so the polynomial--rigidity that forbids reducing the \emph{rational} \#307
to such a hypothesis does not apply here. A Schinzel--type theorem over $\mathbb{Z}[\sqrt2]$ would render
twisted infinitude unconditional, whereas none touches \#307 itself; the twisted problem sits strictly
closer to the reach of current sieve theory than its rational shadow. Concretely the twisted solution set
is \emph{abundant} at accessible scale; thousands of cycles at modest norm (Proposition~\ref{prop:geometry}),
and the count of distinct prime--norm signatures grows roughly in proportion to the norm bound in our census
(e.g.\ $6,12,20$ at bounds $40,80,160$, unit heights saturating); whereas the rational Pythagorean layer
is \emph{empty} below $10^{112}$ (Corollary~\ref{K-cor:emptytest}). The twisted and rational problems thus differ
not in the difficulty of the search but in whether solutions exist to be found at all: for the twist they are
plentiful and the only question is their primality, exactly the Schinzel--shaped residue above. That residue is a
\emph{pair}, not a triple or a Lehmer sequence: fixing $\pi_1$ and bounding the units, the cycle count still
grows with $N(\pi_2)$ (verified: $33,48,60$ for $\pi_1=3+\sqrt2$ at $N(\pi_2)\le60,120,240$), so $\pi_1$'s
primality is free and only the \emph{simultaneous} primality of $\pi_2$ and the forced $\pi_3$ remains. Setting up
the sieve pins the difficulty precisely. Writing $C=u_1\pi_2+u_2\pi_1$, $E'=(\pi_1\pi_2)^\sigma$,
$F=-u_3\pi_1\pi_2$, the forced third prime is $\pi_3=(\bar CF+E'\bar F)/\det$ with $\det=N(C)-N(E')$ a binary
quadratic form in $(c,d)=\pi_2$. The first condition is the clean form $N(\pi_2)=c^2-2d^2$; the second, however,
is \emph{not} a second form but the rational condition $N(\pi_3)=N(G)/\det^2$ ($G=\bar CF+E'\bar F$ quadratic,
$N(G)$ quartic, $\gcd(N(G),\det)=1$), evaluated on the integral locus $\det\mid G$, and $N(\pi_3)$ is unbounded,
spiking as $\det\to0$ (values to ${\sim}10^{11}$ already at $N(\pi_2)\le2\times10^4$). On this \emph{frozen}
$(\pi_1,\text{units})$ slice the question is primality along the folded conic bundle, and the slice is thin: its
unimodular fibres $\det=\pm1$ are \emph{empty} by a genus obstruction (there $\det=\pm1$ would need
$x^2-2y^2\in\{-55,-43\}$, both non--residues). But the thinness is an artefact of freezing exactly the parameters
that carry the abundance. The \emph{full} twisted variety $D(a)=a^\sigma$ ($\pi_1,\pi_2,\pi_3$ and their units all
free) is a cubic threefold whose solution count grows with positive exponent; a positive--density attack must run
there, not on a fixed fibre, so the fixed--fibre conic sieve is a closed dead end. The honest placement is then: the
reachable object is the full threefold at the Heath--Brown $x^3+2y^3$ frontier, compounded by a \emph{coupled} second
primality ($N(\pi_2)$ and $N(\pi_3)$ simultaneously prime); the two--sparse--conditions/parity barrier where the
circle method loses minor--arc cancellation. This is strictly below the exponential/Lehmer tier of the immune
families (Remark~\ref{A-rem:lehmer}, where no sieve reaches at all), but a generation beyond present technology; it is
where genuinely new analytic input for the twist would first have to bite.
\end{remark}

The natural attempt to \emph{prove} that infinitude (forcing one prime to be the conjugate of another
so that a single free prime remains, whence Dirichlet would suffice) is blocked, and the block is a
clean structural law.

\begin{proposition}[No conjugate prime pair in a twisted cycle]\label{prop:noconj}
Let $D$ be a Leibniz derivative with unit prime--values over a quadratic ring with automorphism $\sigma$,
and $a$ a squarefree twisted cycle $D(a)=a^\sigma$. Then no two prime factors of $a$ are
$\sigma$--conjugate. Equivalently a twisted cycle's support is disjoint from its conjugate; it is never
conjugation--stable, which is exactly why the exhibited $\mathbb{Z}[i]$ cycles are asymmetric, coming in
conjugate pairs of \emph{distinct} supports rather than on one symmetric support. This \emph{separates}
two problems easily conflated: the twisted equality $D(a)=a^\sigma$ treated here, and the symmetric
target of Remark~\ref{rem:descent}, which is a genuine two--cycle $D(m)=m'$ ($m\ne m'$, both
conjugation--fixed) of the \emph{descended} rational operator; not a twisted equality. The present
proposition bars the first from the symmetric locus; it does not bear on the second, which remains
genuinely open (the exhaustive emptiness there, Remark~\ref{rem:descent}, is unexplained by it).
\end{proposition}
\begin{proof}
Suppose a split prime $\pi$ and its distinct conjugate $\pi^\sigma$ both divide $a$. Reducing
$D(a)=\sum_i u_i\,(a/\pi_i)$ modulo $\pi$ kills every term but the one at $\pi_i=\pi$, so
$D(a)\equiv u_\pi\,(a/\pi)\not\equiv0\pmod\pi$ ($u_\pi$ a unit and $\pi\nmid a/\pi$ by squarefreeness).
Yet $a^\sigma=\prod_i\pi_i^\sigma$ carries the factor $(\pi^\sigma)^\sigma=\pi$, so
$a^\sigma\equiv0\pmod\pi$; hence $D(a)\ne a^\sigma$. An exhaustive $\mathbb{Z}[\sqrt2]$ search confirms
no cycle carries a conjugate pair.
\end{proof}

We codify what is theorem and what remains open in this circle.

\begin{theorem}[Existence of twisted derivative two--cycles]\label{thm:existence}
Leibniz derivatives with unit prime--values admit genuine two--cycles over $\mathbb{Z}[i]$,
$\mathbb{Z}[\sqrt{-2}]$, $\mathbb{Z}[\omega]$, and $\mathbb{Z}[\sqrt2]$. The proof is by exhibited
certificate, each verified exactly by independent recomputation: the smallest instances are, in order of
norm, $\bigl(\sqrt2(5-\sqrt2),\,7\bigr)$ over $\mathbb{Z}[\sqrt2]$ at norms $(46,49)$; the Eisenstein
cycle at norm $2{,}964$; the sign--twisted cycle over $\mathbb{Z}[\sqrt{-2}]$ at norm $11{,}462$, and the
cycle over $\mathbb{Z}[i]$ at norm $16{,}642$; together with the rational pair
$D(3707)=1547$, $D(1547)=3707$ over $\mathbb{Z}[\sqrt2]$.
\end{theorem}

\begin{proposition}[Geometry of the solution set]\label{prop:geometry}
Over $\mathbb{Z}[\sqrt2]$, consider the Galois--folded system
$u_1\pi_2\pi_3+u_2\pi_1\pi_3+u_3\pi_1\pi_2=(\pi_1\pi_2\pi_3)^\sigma$ of
Remark~\textup{\ref{rem:folding}}. \textup{(i)} For \emph{every} triple of split primes the system is
solvable with the $u_i$ on the real unit hyperbola $\{|st|=1\}$: solving the second embedding for $u_2$
and letting the third coordinate tend to zero, the first--embedding residual sweeps $\pm\infty$ while the
remaining terms stay bounded, so a root exists by continuity; the real solution set is a nonempty curve
\textup(verified numerically on $200$ random triples, all solvable\textup). The volume of the solution
variety never vanishes. \textup{(ii)} By Proposition~\textup{\ref{prop:sunit}} each fixed support carries
finitely many integral solutions, so infinitude can only come from growing supports. \textup{(iii)} The
measured integral hit rates are $11$ of $455$ prime triples in the smallest box, and $30{,}204$ cycles of
norm $\le2\times10^7$ at twist height $\le2$. \textup{(iv)} What remains open is precisely whether the
discrete unit lattice meets the real curve in integer points for infinitely many supports,
inhomogeneous Diophantine approximation on exponentially parametrised curves. Two shortcuts are closed:
bounded--unit polynomial families are obstructed (the rational component of the return forces unit
coefficients to track growing prime coordinates, which bounded units cannot), and even the cleanest
forward family ($D\bigl(\sqrt2(x-\sqrt2)\bigr)=x+2$ \emph{identically in} $x$, supplying a candidate
cycle whenever $x^2-2$ and $x+2$ are prime) has sporadic returns: the closing two--term unit equation,
effectively finite for each $x$, is solvable at $x=5$ and provably insolvable at $x=15$ and $x=21$.
Those failures are, however, exactly the rigid case: a \emph{prime} return $x\pm2$ leaves the closing
system no slack. The identity has four unit branches $D\bigl(\sqrt2(x\pm\sqrt2)\bigr)=x\pm2$ (independent
signs, twists $1$ and $\pm\varepsilon^{\pm1}$), and when $x\pm2$ is composite squarefree the closing
system gains one unit unknown per prime class against the same two equations, and does close:
$x=107,121,387,1693$ yield cycles with rational partners $b=105,119,385,1695$, each verified exactly.
Infinitude along this family is thereby reduced to Bunyakovsky for $x^2-2$ together with recurring
composite--return closure: the softest form the gap takes anywhere in this paper. Even the primality
half of that reduction can be removed. By the half--dimensional sieve
\cite{Iwaniec1978,LemkeOliver2012} (applied to $(2y+1)^2-2=4y^2+4y-1$ to keep $x$ odd), $x^2-2$ is a
$P_2$ infinitely often; every prime factor of $x^2-2$ has $2$ as a quadratic residue and therefore
splits in $\mathbb{Z}[\sqrt2]$, so $\sqrt2(x-\sqrt2)$ is infinitely often squarefree and all--split
with $\omega\in\{2,3\}$. The candidate supports thus exist \emph{unconditionally}; all that remains
open on this route is whether the unit gates open along that supply infinitely often.
Exactness migrates; it does not disappear.
\end{proposition}

\begin{proposition}[No vanishing derivatives, no transport, and the price of exactness]\label{prop:novanish}
\textup{(i)} No squarefree $z$ with $\omega(z)\ge1$ has $D(z)=0$, for any twisted derivative over any of
the rings considered (indeed over any UFD): modulo a prime $\pi\mid z$ every cofactor term but one
vanishes, leaving $D(z)\equiv u_\pi\,z/\pi\not\equiv0\pmod\pi$. The zero--derivative counters of every
census above are therefore theorems, not observations. \textup{(ii)} Consequently cycles cannot be
transported multiplicatively: for squarefree $c$ coprime to $ab$, Leibniz gives $D(ca)=D(c)\,a+c\,D(a)$,
so $(ca,cb)$ is a cycle exactly when $D(c)=0$; impossible. Cycles do not breed; each one requires
fresh exactness. \textup{(iii)} The two cycle equations can be bought by construction, at a conserved
price. Buying \emph{neither} leaves the generic stratum: two exact unit conditions ($30{,}204$
specimens). Buying \emph{one} is free in two dual ways; forward, $\pi\coloneqq q-\varepsilon^m\sqrt2$
makes $D(\sqrt2\,\pi)=q$ an identity (the four--branch family above is the case $m=\pm1$); return,
$\pi\coloneqq(\rho+\varepsilon^{-1}\rho^\sigma)/\sqrt2=(c+2d)-d\sqrt2$ for $\rho=c+d\sqrt2$ closes
$D(b)=\sqrt2\,\pi$ identically for every associate $b$ of $\rho\rho^\sigma$, and the residue is one
primality, $|N(\pi)|$, plus one exact unit gate, measured at $2$ of $50$ pairs with both primalities in
the smallest box. The two hits re--derive, with no search, precisely the two census minimal--shape cycles
with prime rational partner: norms $(46,49)$ and $(206,49)$. Buying \emph{both} collapses the parameters
to unit--indexed quadratics in $q$ with a square--discriminant condition; the resulting $q=7$ line
carries the known cycles at $m=1$ and $m=-2$ (discriminants $25$ and $81$) together with one Galois
image, and every other prime--norm rung out to $|m|\le14$ fails its return, verified exactly. Each
purchased identity is repaid as a sparser arithmetic condition elsewhere: the hardness of exactness is
conserved, never destroyed.
\end{proposition}

\begin{remark}[The two faces of the gap: unit rank dichotomises the obstruction]\label{rem:dichotomy}
The invariant that governs abundance in Remark~\ref{rem:orderable} (the unit rank
$r=r_1+r_2-1$ of Dirichlet) also governs the \emph{type} of the residual open problem, and the two
types are mutually exclusive. When $r=0$ (the rings $\mathbb{Z}$, $\mathbb{Z}[i]$,
$\mathbb{Z}[\sqrt{-2}]$, $\mathbb{Z}[\omega]$) the unit group is finite, so a given pair of supports
admits only $|\mathcal{O}^\times|$ twist assignments and hence finitely many candidate cycles: the
solution locus is a set of \emph{isolated points}, zero--dimensional, and infinitude can arise only from
infinitely many \emph{distinct} supports independently passing a closing test; a
Bateman--Horn/Schinzel--type recurrence over prime tuples, with no continuous family to deform along.
The constituent primalities are individually classical (each support prime is a value of a binary
norm--form, so by Fermat--Dirichlet (and, for $\mathbb{Z}[i]$, by the Green--Sawhney theorem with
\emph{both} coordinates prime) there is no shortage of supports); it is their \emph{simultaneous}
closure that is not classical. This was tested directly: the minimal shape $a=(1+i)\pi$ over
$\mathbb{Z}[i]$ yields \emph{no} cycle for any Gaussian prime $\pi$ of norm $\le4000$ under any of the
sixteen twists, and the sixteen genuine cycles are irregular mixtures of inert and split classes
($\omega=3,4,5$) with no one--parameter family among them; the signature of isolated points, not a
curve. When $r\ge1$ (every real quadratic ring, e.g.\ $\mathbb{Z}[\sqrt2]$),
Proposition~\ref{prop:geometry} makes the real solution locus a positive--dimensional curve in each
support, but the unit lattice is an infinite rank--$r$ discrete set, and infinitude requires it to meet
that curve in \emph{integer} points infinitely often; inhomogeneous Diophantine approximation on an
exponentially parametrised curve. No unit rank delivers both a classical gate and a continuous family:
rank zero gives the former and denies the latter, positive rank the reverse. This is why no single ring
lets the infinitude be \emph{finished} by present methods; a structural account of the obstruction's
universality, not a barrier theorem. It also locates $\mathbb{Z}$ exactly: the unique ring carrying
\emph{both} the rank--zero sporadicity \emph{and}, by formal reality, the mass barrier; the worst
cell of the trichotomy of Corollary~\ref{cor:ff307}, which is precisely why \#307 is the hard
representative of its family.
\end{remark}

\section{Two closures on the existence route}\label{sec:exclosure}

Atom $\mathrm{A}10$ asks for a prime term of the ternary recurrence of
Remark~\textup{\ref{A-rem:lehmer}}. Before attacking it, two families of approach can be removed:
the relaxations that would give the construction more freedom, and the algebraic methods that
would avoid the archimedean place. Both are closed outright.

\begin{proposition}[No free slot]\label{prop:nofreeslot}
Let $P,Q$ be finite sets of primes with $\sigma(P)\sigma(Q)=1$, and put $a=\prod P$, $b=\prod Q$.
Then $b=a'$ and $a=b'$ exactly. In particular $Q$ is \emph{determined} by $P$: it is the set of
prime factors of $a'$, and the existence route carries exactly one parameter, the set $P$.

\textup{Formalised.} \texttt{Erdos307.Q\_determined}: given $a'=b$, the set $Q$ is exactly the set of
prime factors of $a'$. The one--prime normal form of Proposition~\ref{K-prop:oneprime} is
\texttt{oneprime\_forced}, \texttt{oneprime\_cycle}, \texttt{oneprime\_converse}
\textup{(}which closes the equivalence\textup{)},
\texttt{oneprime\_mass\_identity} and \texttt{oneprime\_mass\_bound}, the last two giving
part~\textup{(b)}: the mass identity holds for
\emph{every} witness, prime or composite, so the exactness layer alone carries the barrier. The slot
calculus of Proposition~\ref{C-prop:slot} is \texttt{slot\_layers} and \texttt{slot\_recovery}, the
filter of Corollary~\ref{C-cor:qr} is \texttt{qr\_filter}, resting on \texttt{cofactor\_survival},
the congruence $N'+2N\equiv N/\ell\pmod\ell$, which is proved rather than assumed (stating \texttt{qr\_filter} with that arithmetic as a hypothesis on abstract residues would leave it contentless); Lemma~\ref{C-lem:symbolfact} is \texttt{Frame.csum\_eq\_mul\_recipSum} for its first half
and \texttt{symbolfact} for the Jacobi factorisation; and the reduction mechanism behind
Proposition~\ref{prop:noconj}, Rigidity and Lemma~\ref{A-lem:anatomy} alike is isolated once as
\texttt{dvd\_sum\_erase\_iff} \textup{(\texttt{lean/Erdos307/NormalForm.lean})}.
\end{proposition}
\begin{proof}
$\sigma(a)=a'/a$ and $\sigma(b)=b'/b$, both in lowest terms since $\gcd(m,m')=1$ for squarefree
$m$. Then $\sigma(a)\sigma(b)=1$ reads $a'b'=ab$. From $\gcd(a',a)=1$ and $a'\mid ab$ we get
$a'\mid b$, say $b=a'k$; symmetrically $b'\mid a$, say $a=b'j$. Substituting,
$a'b'=ab=b'j\cdot a'k$, so $jk=1$ and $j=k=1$.
\end{proof}

\begin{corollary}[The multi-slot relaxations are vacuous]\label{cor:noslack}
There is no version of the construction in which two or more primes of $Q$ may be chosen freely.
In particular one cannot fix a base $B\subset Q$ and solve $1/q_1+1/q_2=r/s$ for the remaining two
primes by ranging over the $\tau(s^2)$ factorisations $(rq_1-s)(rq_2-s)=s^2$: by
Proposition~\ref{prop:nofreeslot} the pair $(q_1,q_2)$ is already forced to be the corresponding
part of the factorisation of $a'$. The exponential parametrisation of
Remark~\textup{\ref{A-rem:lehmer}} is therefore intrinsic, not an artefact of the one-free-prime
ansatz, and no amount of extra slots converts it into a Bunyakovsky one. What does dispose of the
exponential growth is not extra slots but Proposition~\ref{K-prop:tailbound}, which bounds the tail
prime outright and so replaces the orbit by a finite divisor search.
\end{corollary}

The second closure is not new here, but it belongs in this list because it is what forces the whole
existence route to be analytic. Theorem~\ref{thm:ff} of Section~\ref{sec:ff} shows that over
$\mathbb{F}_q[t]$ the arithmetic derivative satisfies $\deg f'\le\deg f-1$, hence has no two--cycles
and no nonzero fixed points, and Corollary~\ref{cor:ff307} concludes that the function--field
Erd\H{o}s--Barbeau equation is insoluble outright. \textup{(An independent brute-force recheck,
$8192$ squarefree monics over $\mathbb{F}_2$ to degree $13$ and $6561$ over $\mathbb{F}_3$ to
degree $8$, $0$ violations and $0$ two--cycles: \texttt{code/ff\_nocycle.rs}.)} What is added here
is not that statement but its consequence for \emph{method}, recorded next.

\emph{Why the existence route cannot be algebraic.} The mechanism of Theorem~\ref{thm:ff} is the ultrametric inequality. The cycle equation is
$\sigma(a)\sigma(b)=1$, so one of the two masses must reach $1$; over $K[t]$ the mass satisfies
$|\sigma(f)|=\max_i|1/\pi_i|<1$ by ultrametricity and can never do so, whereas over $\mathbb{Z}$ the
archimedean absolute value permits $\sum_{p\mid n}1/p>1$. So \#307 is possible at all only because
$|\cdot|_\infty$ is not ultrametric.

This has a consequence for method. Any argument for \emph{existence} that is purely algebraic, in
the sense that it applies verbatim with $\mathbb{Z}$ replaced by $K[t]$, would prove a statement
that Theorem~\ref{thm:ff} refutes; hence no such argument exists. Together with
Proposition~\ref{K-prop:nopoly} (no polynomial parametrisation) and Proposition~\ref{C-prop:noplace}
(no place of $\mathbb{Q}$ carries the descent) this closes the algebraic and geometric toolkit on
\emph{both} routes: the negative route has no place to descend at, and the positive route has no
place to construct at. What remains for $\mathrm{A}10$ is archimedean and analytic, which is
precisely the regime in which lower bounds for primes in sparse sequences are unavailable. That is
a sharper statement of the obstruction than Remark~\ref{A-rem:lehmer} gives, and it is the reason the
present work does not resolve the existence direction.


A third closure was claimed here in place of what follows, and it was wrong. The $34$ immune
families of Remark~\textup{\ref{A-rem:campaign}} are precisely the tail families the congruence
campaign cannot kill, so each is a live candidate and the obvious move is to test them one at a
time. That move \emph{is} available, and the argument that said otherwise is worth setting out,
because its error is a natural one.

\begin{remark}[Why the immune families looked untestable, and why they are not]\label{rem:untestable}
Deciding whether the tail family $S\cup\{q\}$ contains a Pythagorean pair is, by
Remark~\textup{\ref{A-rem:lehmer}}, deciding whether $B_Sx^2-A_Sy^2=-4D_S^2$ has an integral solution
with $q=(x^2-D_S)/A_S$ prime, and two routes to that decision present themselves.
\begin{enumerate}[label=\textup{(\roman*)},leftmargin=2.2em]
\item \emph{Global route.} Enumerating \emph{all} integral points on the conic requires the
fundamental unit, equivalently the class group, of the real quadratic order of discriminant
$4A_SB_S$. On the first immune base that discriminant has $225$ digits and is a nonsquare, so the
subexponential index-calculus cost
$L_\Delta(1/2)=\exp\bigl(\sqrt{\log\Delta\,\log\log\Delta}\bigr)$ evaluates to $10^{24.7}$
operations. This remains true, and it is what blocks the squarefree--density argument discussed
above, which genuinely needs the Pell orbit modulo $p^2$.
\item \emph{Parametric route.} Since $A_S$ is prime and $(D_S\mid A_S)=+1$, a square root
$x_0^2\equiv D_S\pmod{A_S}$ exists and the solutions are parametrised by $x=x_0+tA_S$ with
$q(t)=\bigl((x_0+tA_S)^2-D_S\bigr)/A_S$. A solution additionally needs $B_Sq(t)+D_S$ to be a square,
a quantity of $223$ digits already at $t=0$, so the heuristic hit rate is $10^{-112}$ per trial and
the parameter $t$ appears to range without bound.
\end{enumerate}
It does not. The earlier reading of (ii) overlooked the bound this paper had already proved:
Proposition~\ref{K-prop:tailbound} forces $q\le tD_S$ and hence $x<A_S$
(Proposition~\ref{A-prop:lvl60factor}), so $t=0$ and the two square roots $x_0$ and $A_S-x_0$ exhaust
the parametrisation. The decision is therefore a pair of primality tests, not an unbounded sweep,
and Proposition~\ref{A-prop:immunedecide} carries it out: all $34$ immune families are empty. The
global route (i) is not needed because the decision never asks for all integral points on the conic,
only for those in a box the tail bound has already made small. The error was to treat a
parametrisation as a search space without asking what the accompanying inequality did to it, and it
cost the paper a false impossibility claim.
\end{remark}

One hope survives all of this on its face: a lower bound for \emph{some term of some sequence in a
family} is formally weaker than one for a single sequence, and Proposition~\ref{K-prop:close59}
supplies $49{,}961$ of them. It fails, and it fails for a reason no enlargement of the family can
repair.

\begin{proposition}[Multiplicity cannot help]\label{prop:multiplicityfail}
Every known lower-bound method for primes in a set requires that set to have counting function at
least $X^{\theta}$ for some fixed $\theta>0$. The union over all $49{,}961$ bases of the terms of
their sequences below $X$ has counting function $C\log X$ with $C=49{,}961/\log\varepsilon$, since
each sequence grows like $\varepsilon^n$ and so contributes only $O(\log X)$ terms. Hence
\[ \frac{\log(\text{union below }X)}{\log X}=\frac{\log\bigl(C\log X\bigr)}{\log X}\;\longrightarrow\;0, \]
and the family size $C$ is powerless because it sits inside a logarithm.
\end{proposition}
\begin{proof}[Measurement]
On an immune base $\varepsilon$ has $224$ digits, so at $X=10^{500}$ each sequence contributes
$2.2$ terms and the whole family contributes $10^{5.05}$, a density exponent of $0.0101$. Reaching
$\theta=\tfrac12$ at that $X$ would need $10^{249.7}$ bases against the $10^{4.7}$ available, a
shortfall of $10^{245}$. Even the absolute ceiling on the family at \emph{every} level combined,
all $2^{139}=10^{41.8}$ subsets of the primes below $800$, falls short by $10^{207.8}$
\textup{(\texttt{code/a10\_multiplicity.gp})}.
\end{proof}

Each of the four blocks was then attacked on its own terms. None fell; each is now located more
precisely, and two further routes are closed.

\begin{proposition}[The number--field barrier transfer is \textsc{false}]\label{prop:fieldoptimal}
Theorem~\ref{thm:ff} kills the algebraic route by ultrametricity, which suggests replacing $K[t]$ by
a number ring $\mathcal{O}_K$, where the absolute value is archimedean and the obstruction
disappears. It does disappear, so it is not the case that for every number field $K$ the least $\prod N(\pi)$ over prime ideals with $\sum1/N(\pi)\ge2$ is
at least that of $\mathbb{Q}$. That assertion is false. What holds is the following.
\begin{enumerate}[label=\textup{(\roman*)},leftmargin=2.2em]
\item \emph{The true infimum is $16$.} In the biquadratic field
$K=\mathbb{Q}(\sqrt{17},\sqrt{33})$, of discriminant $314{,}721$, the prime $2$ splits completely:
there are four prime ideals of norm $2$. They carry mass $4\cdot\tfrac12=2$ exactly, and their
product of norms is $2^4=16$. Conversely $2$ is the least possible ideal norm, so mass $2$ needs at
least four ideals of norm $2$ and any admissible family has product at least $16$; a field can carry
at most $[K:\mathbb{Q}]$ ideals of norm $2$, so degree $4$ is also minimal. Hence $16$ is the exact
infimum over all number fields, against $\prod_{i\le59}p_i>8.77\times10^{112}$, a factor $10^{112}$ apart.
\item \emph{Where the argument broke.} The Chebotarev computation is correct as a statement about
\emph{density}: in a degree-$d$ Galois field the completely split primes have density $1/d$, so they
carry mass $d\cdot(1/d)\log\log Y=\log\log Y$, exactly as over $\mathbb{Q}$, while each costs
$d\log p$, and inert primes contribute $1/p^{d}$, that is, nothing. The error was to read a
statement about the tail as a statement about the minimum. The barrier consumes only \emph{finitely
many} primes, and for any finite set of small primes one may choose $K$ in which all of them split;
Chebotarev constrains how often splitting happens, not which particular primes split.
\item \emph{What survives, and what closes the route instead.} For a \emph{fixed} field the
measurement stands: $\mathbb{Q}$ needs $59$ primes and $10^{112.9}$, whereas $\mathbb{Q}(i)$ needs
$226$ prime ideals and $10^{735.1}$ \textup{(\texttt{code/a10\_fieldbarrier.gp})}. What does not
survive is the use of this proposition as a \emph{block}. The route is nevertheless closed, and by a
sounder argument, which is Proposition~\ref{prop:distinct} below: the ideal--theoretic \#307 is
\emph{trivially soluble}, so it carries no information about the rational problem.
\end{enumerate}
\end{proposition}

\begin{proposition}[The difficulty of \#307 is distinctness, not mass]\label{prop:distinct}
Allow $P$ and $Q$ to be finite \emph{multisets} of primes rather than sets. Then the \#307 equation
is soluble with four primes in total:
$\bigl(\tfrac12+\tfrac12\bigr)\bigl(\tfrac12+\tfrac12\bigr)=1$. With $P,Q$ genuine sets the same
equation forces $|P\cup Q|\ge59$ and $\prod>10^{112.9}$ by Theorem~\ref{K-thm:barrier}. The barrier is
therefore a statement about the injectivity of $p\mapsto1/p$ on the primes, not about the arithmetic
of the mass condition.

Two consequences, and the second is the operative one.

\textup{(i)} \emph{The ideal--theoretic problem is trivial.} In $K=\mathbb{Q}(\sqrt{17},\sqrt{33})$
the prime $2$ splits completely into four pairwise coprime prime ideals of norm $2$; taking
$P=\{\pi_1,\pi_2\}$ and $Q=\{\pi_3,\pi_4\}$ gives $\sigma(P)=\sigma(Q)=1$ and $|P\cup Q|=4$
\textup{(\texttt{code/ideal307.gp})}. Nor does the derivative transfer: $\pi_1+\pi_2$ is the unit
ideal, so $a'$ collapses. This is the correct reason the number--field route is closed, and it
explains why refuting the previous proposition reopened nothing: the transferred problem was never
the same problem.

\textup{(ii)} \emph{A no--go for transfers generally.} Any setting admitting two distinct objects of
the same norm $n\ge2$ on each side solves the mass equation at size four, whatever the setting is.
So a dictionary permitting repeated denominators \emph{dissolves} \#307 rather than relocating its
difficulty, and can carry no information about it. This prunes a class of attacks in one line, and
it is the test any future transfer should face first. Corollary~\ref{K-cor:coprime60} passes it only
because pairwise--coprime integers still have distinct least prime factors, which is distinctness
re-entering through the back door.

\textup{Formalised.} \texttt{Erdos307.multiset\_solution} is the four--element witness,
\texttt{distinctness\_carries\_the\_barrier} the contrast against
\texttt{card\_ge\_59\_of\_recipSum\_ge\_two}, stated as a conjunction so the two halves cannot drift
apart, and \texttt{transfer\_dissolves} the general no--go
\textup{(\texttt{lean/Erdos307/Distinctness.lean})}.
\end{proposition}

\emph{The three other blocks, attacked.} \textup{(i) Against Proposition~\ref{prop:nofreeslot}.} The rigidity $\prod Q=a'$ would dissolve if
disjointness were an assumption one could relax. It is not: $P\cap Q=\varnothing$ is
\emph{derived} in Theorem~\ref{K-thm:structure} by the Israel mechanism, $r\in P\cap Q$ forcing
$v_r(st)=-2\neq0$. There is nothing to relax.

\textup{(ii) The block that fell.} A block was claimed here on the ground that deciding an immune
family needs integral points on the conic $B_Sx^2-A_Sy^2=-4D_S^2$, hence the fundamental unit of a
$225$--digit discriminant. The objection to it was that conics obey Hasse--Minkowski, so with the
local solubility of Remark~\ref{A-rem:tier60} rational solubility is unconditional; the reply was that
what is missing is integrality, not solubility. Both sides of that exchange were beside the point.
Integral points are not needed either, because Proposition~\ref{K-prop:tailbound} confines the solution
to a box containing at most two candidates, and Proposition~\ref{A-prop:immunedecide} simply tests
them. Remark~\ref{rem:untestable} records the error. The lesson is worth more than the block was: an
obstruction stated in terms of a general algorithmic problem is only as good as the argument that the
problem instance is general, and here it was not.

\textup{(iii) Against Proposition~\ref{prop:multiplicityfail}.} The template that exploits a family
without density is Heath-Brown's: Artin's conjecture holds for all but at most two primes, with no
means of saying which. It does not transfer. That argument needs each failure to impose a
\emph{detectable} structure, which several failures then cannot share. Here a failing family would
be one all of whose $q_n$ are composite, and the only known mechanism for that is a covering system,
which Proposition~\ref{A-prop:localcomplete} proves does not exist. A failure would therefore be
structureless, and structureless failures cannot be played against one another.


All of the above bounds the negative direction. The positive direction admits a quantitative
shadow that had not been measured: a solution is exactly $r(a):=\sigma(a)\sigma(a')=1$, so one may
ask how close $r$ actually comes to $1$. The answer is instructive, and it disposes of the usual
way such heuristics are calibrated.

\begin{proposition}[The near-miss spectrum, and why the heuristic cannot be calibrated]\label{prop:nearmiss}
Over all $4{,}385{,}727$ integers $a\le10^7$ with $a$ and $a'$ both squarefree, the maximum of
$r(a)$ is $0.5535$, attained at $a=3{,}972{,}254$; the frontier advances by decade as
$0.333,\,0.345,\,0.502,\,0.502,\,0.536,\,0.554$, and at $10^{8}$ it reaches $0.586076$ at
$a=23{,}667{,}105$ \textup{(\texttt{code/sweep1e8.rs}, which reproduces every earlier decade
independently)}. This last decade was run as a test of Proposition~\ref{prop:frontier}: the bound
$R(10^{8})\le0.67270$ and the range $[0.55,0.62]$ suggested by the earlier decades were both
recorded before the computation, and both hold. In particular $r\ge1$ never occurs, the defect
$|a''-a|$ is of size $a^{0.9}$ or larger in all but three cases, and the smallest relative defect is
$0.4465$. Consequently the constant in the heuristic count of two-cycles, $f(1)\log X$ with $f$ the
density of $r$ at $1$, is \emph{not measurable}, and the obstruction is two--sided.

By Theorem~\textup{\ref{K-thm:barrier}} the region $r\ge1$ lies above $10^{112.9}$. The window
immediately to the \emph{left} of $1$ is barred as well, by AM--GM rather than by the barrier
hypothesis: $\sigma(a)+\sigma(a')\ge2\sqrt{\sigma(a)\sigma(a')}=2\sqrt r$, and
$2\sqrt r>T_{58}=1.99874\ldots$ as soon as $r>(T_{58}/2)^2=0.998740\ldots$, whence $|P\cup Q|\ge59$
and the product again exceeds $\Pi_{59}>8.77\times10^{112}$. So the whole neighbourhood
$r\in(0.998740,\infty)$ of the value $1$ lies above the barrier, not merely the half--line $r\ge1$,
and no sweep reaches any of it \textup{(\texttt{code/a10\_nearmiss.rs},
\texttt{code/a10\_density1.rs})}. Barring only $r\ge1$ would leave the left window formally open and the unmeasurability argument one--sided.
\end{proposition}

\begin{proposition}[An effective approximation theorem, with a certified witness]\label{prop:effapprox}
Bado's Theorem~10.2 \cite{Bado2026} asserts that for every $Y$ and every $\varepsilon>0$ there are
disjoint finite sets of primes $P,Q$, all greater than $Y$, with $0\le\sigma(P)\sigma(Q)-1<
\varepsilon$. The proof is a greedy accumulation and is ineffective. Made effective it is expensive:
to land in $[T,T+\eta)$ the accumulation must start above $\approx1/\eta$, and Mertens gives
$\sum_{Y_1<p\le W}1/p\approx\log\log W-\log\log Y_1=T$, so $W=Y_1^{e^{T}}$; two stages with
$T\approx1$ put the largest prime and the cardinality both at $\approx\varepsilon^{-e^{2}}=
\varepsilon^{-7.389\ldots}$.

A Newton ladder does the same work with a handful of primes. Take $Q=\{2,3,5\}$, so
$\sigma(Q)=31/30$ and $P$ must approach $30/31$; from a pool, repeat $p\leftarrow$ the least prime
$\ge1/d$ and $d\leftarrow d-1/p$, where $d$ is the running deficit. Each rung squares it,
$d_{\mathrm{new}}=d-1/p\approx t\,d^{2}$ with $t$ the local prime gap, so the largest prime is
$\approx\varepsilon^{-1/2}$ and the number of rungs is $O(\log\log(1/\varepsilon))$. Taking the pool
to be the primes $7\le p\le271$ and eleven rungs gives $|P|=66$, $|Q|=3$, largest prime of $2{,}010$
digits, and, in exact rational arithmetic,
\[ \bigl|\sigma(P)\sigma(Q)-1\bigr|\;=\;4.98523\ldots\times10^{-4016}, \]
a rational whose numerator has $129$ digits and whose denominator has $4{,}145$. Every rung prime is
\emph{proved} prime by \textup{APR--CL}, not merely a pseudoprime
\textup{(\texttt{code/nearmiss\_effective.gp})}. The certificate is tiered, and a reader may stop
where they like:
\[
\begin{array}{l|ccc}
\text{rungs} & 9 & 10 & 11\\\hline
\text{largest prime (digits)} & 505 & 1{,}007 & 2{,}010\\
\bigl|\sigma(P)\sigma(Q)-1\bigr| & 9.57\times10^{-1007} & 2\times10^{-2010} &
  4.99\times10^{-4016}\\
\text{cumulative certification} & <10\text{ s} & \approx110\text{ s} & \approx27\text{ min}
\end{array}
\]
The statement is machine--checked at the three--rung tier, where the largest prime is
$2{,}348{,}039{,}453$ and the defect is $2.1111\ldots\times10^{-18}$: the primality of all $58$
primes, the disjointness of the two sides, the exact value of $\sigma(P)$ as a ratio of a $126$--
and a $129$--digit integer, and the bound $|\sigma(P)\sigma(Q)-1|<10^{-17}$ are all verified in
\texttt{Erdos307.Witness.effapprox\_witness}, on the standard axioms and with no appeal to
kernel--external evaluation. Beyond that tier the primes are certified by APR--CL, which Mathlib
does not carry; the higher rungs are therefore cited to the \textup{PARI/GP} certificate rather
than formalised. $0$ \texttt{sorry} \textup{(\texttt{lean/Erdos307/Witness.lean})}.
\end{proposition}

The gap is not marginal. At $\varepsilon=10^{-13}$ the greedy needs a largest prime near $10^{94}$
and some $10^{92}$ primes; the ladder needs a largest prime near $10^{7}$ and three rungs. The
ladder is exponentially better in the size of the largest prime and doubly exponentially better in
the count, and the reason is visible: greedy convergence is linear in the number of primes, Newton
convergence is quadratic. Truncating gives certificates of any size, five rungs and a $34$--digit
largest prime already give $6.07\times10^{-66}$, and seven rungs give $3.22\times10^{-254}$. The witness above approaches $1$ from below; taking
the last rung's prime just under $1/d$ instead overshoots, so either side is available and the
statement of Theorem~10.2 as printed is met.

What this is not: progress towards a solution. A free pair of disjoint prime sets carries none of
the rigidity an exact hit requires ($\sigma(P)\sigma(Q)=1$ forces $\prod P=(\prod Q)'$ and
$\prod Q=(\prod P)'$, and nothing in the ladder produces that) and by
Proposition~\ref{prop:nearmiss} the size of $|\sigma(P)\sigma(Q)-1|$ does not measure progress
towards the integer identity at all. The value of the construction is that it makes an existence
theorem effective and exhibits a witness anyone can check.

\begin{proposition}[The near--miss frontier obeys a double--logarithmic law]\label{prop:frontier}
For $X\ge2$ put $R(X)=\max\{r(a):a\le X\ \text{squarefree}\}$, where $r(a)=\sigma(a)\sigma(a')$. Let
$\omega_{\max}(X)$ be the largest $k$ with $\Pi_k\le X$, let $\sigma_{\max}(X)=T_{\omega_{\max}(X)}$,
and let $n(X)$ be the largest $n$ with $\Pi_n\le X^{2}\sigma_{\max}(X)$. Then
\[ R(X)\ \le\ \Bigl(\tfrac12\,T_{n(X)}\Bigr)^{2}. \]
\end{proposition}
\begin{proof}
Rigidity gives $\gcd(a,a')=1$, so $U=\operatorname{supp}(a)\sqcup\operatorname{supp}(a')$ has
$|U|=\omega(a)+\omega(a')$ distinct primes and $\prod_{p\in U}p\mid aa'$. By AM--GM,
$\sum_{p\in U}1/p=\sigma(a)+\sigma(a')\ge2\sqrt{\sigma(a)\sigma(a')}=2\sqrt{r}$, while the same sum
is at most $T_{|U|}$ by the extremality of the smallest primes; hence $r\le(T_{|U|}/2)^{2}$. For the
size, $\Pi_{|U|}\le\prod_{p\in U}p\le aa'=a^{2}\sigma(a)\le X^{2}\sigma_{\max}(X)$, so
$|U|\le n(X)$, and $T$ is increasing.
\end{proof}

The bound is unconditional, needs no squarefreeness of $a'$, and is close: against the measured
frontier of Proposition~\ref{prop:nearmiss} it reads
\[
\begin{array}{l|cccccc}
X & 10^{2} & 10^{3} & 10^{4} & 10^{5} & 10^{6} & 10^{7}\\\hline
\text{bound} & 0.4014 & 0.4920 & 0.5296 & 0.5879 & 0.6129 & 0.6342\\
\max r \text{ observed} & 0.333 & 0.345 & 0.502 & 0.502 & 0.536 & 0.554
\end{array}
\]
\textup{(\texttt{code/frontier\_bound.gp})}. Two consequences. First, since
$T_{n(X)}=\log\log\bigl(X^{2}\sigma_{\max}\bigr)+M+o(1)$ by Mertens, the frontier can climb only like
\[ R(X)\ \le\ \tfrac14\bigl(\log\log X^{2}+M+o(1)\bigr)^{2}, \]
a \emph{doubly logarithmic} law: this is why the measured frontier crawls, and why an extrapolation
of its per--decade advance is worthless. Second, the bound reaches $1$ exactly when $T_{n}\ge2$, that
is $n\ge59$, that is $X^{2}\sigma_{\max}\ge\Pi_{59}$, giving $X\ge2.0942\times10^{56}$, the
constant of Theorem~\ref{K-thm:barrier}, recovered here from the frontier side rather than from the
cycle equations. The two derivations are independent and agree to all printed digits.

\begin{corollary}[The cost of a near--miss; the barrier as a function of $r$]\label{cor:cost}
For $r>0$ put $n(r)=\min\{n:T_n\ge2\sqrt r\}$. Every squarefree $a$ with $r(a)\ge r$ has
$|U|\ge n(r)$, hence $aa'\ge\Pi_{n(r)}$ and
\[ \min(a,a')\ \ge\ \sqrt{\Pi_{n(r)}/T_{n(r)}}. \]
\end{corollary}
\begin{proof}
The first display of the proof of Proposition~\ref{prop:frontier} gives $2\sqrt r\le\sum_{p\in U}1/p
\le T_{|U|}$, so $|U|\ge n(r)$; the rest is Theorem~\ref{K-thm:barrier} verbatim with $59$ replaced by
$n(r)$.
\end{proof}

At $r=1$ this returns $n=59$ and $2.0929\times10^{56}$, so Theorem~\ref{K-thm:barrier} is the endpoint
of a one--parameter family rather than an isolated statement. The family is what makes the crawl
quantitative \textup{(\texttt{code/nearmiss\_cost.gp})}:
\[
\begin{array}{l|cccccccc}
r & 0.5 & 0.554 & 0.6 & 0.7 & 0.8 & 0.9 & 0.99 & 1\\\hline
n(r) & 8 & 9 & 11 & 16 & 24 & 37 & 56 & 59\\
\min(a,a')\ge & 2.6\!\cdot\!10^{3} & 1.2\!\cdot\!10^{4} & 3.6\!\cdot\!10^{5} & 4.4\!\cdot\!10^{9} &
 1.2\!\cdot\!10^{17} & 4.3\!\cdot\!10^{30} & 4.7\!\cdot\!10^{52} & 2.1\!\cdot\!10^{56}
\end{array}
\]
Each increment in $r$ is paid for exponentially in the product, which is the same double--logarithmic
law seen from the other side: $r$ enters the bound only through $2\sqrt r$, and $T_n$ grows like
$\log\log p_n$. It also explains the shape of the swept data. The record $r=0.5535$ needs
$\min(a,a')\ge1.2\times10^{4}$ and was found at $a=3{,}972{,}254$; the next decade of $r$ costs four
decades of $a$, and by $r=0.7$ the requirement has passed $10^{9}$, beyond any sweep that has been
run.

The restriction $r\le1$ is not used in the proof, and above $1$ the corollary says how expensive it
is for the second derivative to \emph{exceed} its argument:
\[
\begin{array}{l|cccccc}
r & 1 & 1.03944 & 1.25 & 1.5 & 2 & 9/4\\\hline
n(r) & 59 & 71 & 212 & 940 & 37{,}963 & 361{,}139\\
\min(a,a')\ge & 10^{56.3} & 10^{71.3} & 10^{274.8} & 10^{1579.7} & 10^{98275.7} & 10^{1127769.3}
\end{array}
\]
So $a''\ge2a$ forces a union support of $37{,}963$ primes and $\min(a,a')>10^{98275}$: the second
derivative cannot double its argument except at wholly inaccessible size. The witness of
Theorem~\ref{thm:lyapfalse}, with $r=1.03944$, needs $|U|\ge71$ and $\min\ge10^{71.3}$, and its $142$
digits supply both.

\begin{remark}[One inequality, six barriers]\label{rem:dictionary}
Corollary~\ref{cor:cost} and Propositions~\ref{K-prop:kuniform} and~\ref{K-prop:twoparam} are the same
statement read at different constants. Every barrier in this note is an instance of
\[ \sum_{p\in U}\frac1p\ \ge\ c \quad\Longrightarrow\quad |U|\ \ge\ n(c),\qquad \prod_{p\in U}p\ \ge\
\Pi_{n(c)}, \]
with $n(c)=\min\{n:T_n\ge c\}$, and the arithmetic enters only in producing $c$:
\[
\begin{array}{l|c|c|l}
\text{configuration} & c & n(c) & \text{recorded as}\\\hline
\text{near--miss ratio }r & 2\sqrt r & \text{varies} & \text{Corollary~\ref{cor:cost}}\\
r=1,\ \text{i.e.\ a two--cycle} & 2 & 59 & \text{Theorem~\ref{K-thm:barrier}}\\
k\text{--cycle} & k/\lfloor k/2\rfloor & 59\ (k\text{ even}) & \text{Proposition~\ref{K-prop:kuniform}}\\
\text{three--cycle} & 3 & 361{,}139 & \text{Proposition~\ref{K-prop:kcycles}}\\
\text{all--odd, }\min U\ge3 & 2\ \text{on }p\ge3 & 1{,}412 & \text{Theorem~\ref{K-thm:parity}}\\
a''\ge2a & 2\sqrt2 & 37{,}963 & \text{above}\\
\text{orbit of growth }G\text{ in }L\text{ steps} & \tfrac{L}{\lceil L/2\rceil}G^{1/L} &
 \text{varies} & \text{Proposition~\ref{prop:growth}}
\end{array}
\]
The coincidence $2\sqrt{9/4}=3=c(3)$ is why $r=9/4$ and the three--cycle carry the identical bound
$361{,}139$: they are the same value of $c$. What looked like six separate barriers is one
mass--to--size dictionary, and the only mathematics specific to each row is the derivation of its
$c$.
\end{remark}

\begin{proposition}[The mass window, and the forced core]\label{prop:window}
Let $(a,b)$ be a two--cycle with $|U|=n$ and let $t_n>1$ solve $t+1/t=T_n$. Then
\[ \frac1{t_n}\ \le\ \sigma(a),\,\sigma(b)\ \le\ t_n, \]
and every prime $p$ with $1/p\ge T_{n+1}-2$ lies in $U$.
\end{proposition}
\begin{proof}
The masses are positive with product $1$ and sum at most $T_n$, so each is between the roots of
$z^2-T_nz+1$. If $p\notin U$ then the $n$ primes of $U$ carry at most $T_{n+1}-1/p$, and
$\sigma(U)>2$ forces $1/p<T_{n+1}-2$.
\end{proof}

Both halves are sharp and both are informative at $n=59$. The window is
$[0.952682,\,1.049668]$: \emph{each member's mass is pinned within five per cent of $1$}. And the
forcing puts every prime up to $167$ in $U$, $39$ of the $59$, so two thirds of the support is
determined before any search begins. The window also recovers
Proposition~\ref{prop:split}: the smaller member needs $T_k\ge1/t_n$, which gives $k\ge3$ for
$n\le68$, $k\ge2$ from $n=69$ and $k\ge1$ at $n=1413$, reproducing that proposition's three
thresholds exactly. The forcing decays with the level:
\[
\begin{array}{l|cccccccc}
n & 59 & 60 & 61 & 62 & 63 & 65 & 68 & 72\\\hline
\text{all }p\le & 167 & 103 & 73 & 61 & 47 & 37 & 23 & 19\\
\text{forced} & 39 & 27 & 21 & 18 & 15 & 12 & 9 & 8
\end{array}
\]
\textup{(\texttt{code/forced\_core.gp})}. This accounts for the size of the level--$59$ enumeration.
The $39$ forced primes carry mass $1.91162\ldots$, so the remaining $20$ must supply more than
$0.08838$ from primes above $167$; the twenty smallest such primes reach exactly that, at $277$. The
free part is therefore a choice of $20$ primes from a narrow band, which is why the admissible count
is $49{,}961$ and not astronomical. At level $60$ the forced core has fallen to $27$ and the free
part to $33$, and by level $68$, the top of the window Proposition~\ref{prop:split} says a level
search must cross, only $9$ primes remain forced.

\begin{proposition}[The free band; why exactly one level is finite]\label{prop:band}
Let $U$ be an admissible support with $|U|=n$ and let $q=\max U$. If $T_{n-1}<2$ then
\[ q\ <\ \frac1{2-T_{n-1}}. \]
For $n\le58$ no admissible support exists, since $T_n<2$. For $n=59$ the hypothesis holds,
$T_{58}=1.99874\ldots$, and $q<793.68$, so $q\le787$. For $n\ge60$ it fails, since
$T_{59}=2.00235\ldots>2$, and no bound on $q$ is available.
\end{proposition}
\begin{proof}
The $n-1$ primes of $U$ other than $q$ carry at most $T_{n-1}$, so $\sigma(U)>2$ forces
$1/q>2-T_{n-1}$.
\end{proof}

This is the mechanism of Proposition~\ref{K-prop:levelbarrier} in its sharpest form. Level--by--level
enumeration is finite at $59$ and infinite from $60$, and the switch is the single fact
\[ T_{58}\ <\ 2\ <\ T_{59}. \]
Below $59$ there is not enough mass to reach $2$ at all; at $59$ every prime is capped at $787$; at
$60$ the first $59$ primes already carry mass $2.00235>2$, so appending \emph{any} prime leaves an
admissible support and the count is infinite at once. There is nothing gradual about the transition,
and it happens at one level.

Together with Proposition~\ref{prop:window} this closes level $59$ from both sides. The support
contains every one of the $39$ primes $\le167$ and is contained in the $138$ primes $\le787$, so the
free part is a choice of $20$ primes from the $99$ candidates in $(167,787]$. Unconstrained that is
$\binom{99}{20}=4.288\times10^{20}$ choices; the mass condition $\sigma(U)>2$ cuts it to $49{,}961$,
a reduction by a factor $8.6\times10^{15}$ \textup{(\texttt{code/free\_band.gp})}. The level--$59$
computation of Proposition~\ref{K-prop:close59} is therefore not a brute force over an unbounded space
but an exact traversal of a set whose two ends are both pinned by the same inequality.

\begin{corollary}[The level--$59$ count, from the two pins]\label{cor:count59}
The admissible supports at level $59$ are exactly the sets
\[ \{p:p\le167\}\ \cup\ S, \qquad S\subset\{p:167<p\le787\},\quad |S|=20,\quad
   \sum_{p\in S}\frac1p\ >\ 2-T_{39}\,=\,0.0883759429\ldots, \]
and there are $49{,}961$ of them.
\end{corollary}
\begin{proof}
Membership of the core is Proposition~\ref{prop:window}, the ceiling is Proposition~\ref{prop:band},
and the displayed inequality is $\sigma(U)>2$ after subtracting the forced mass
$T_{39}=1.9116240570\ldots$. The count is a $20$--from--$99$ knapsack, evaluated in
\texttt{code/level59\_count.rs}.
\end{proof}

This is not a closed formula, subset--sum counts do not have one, but it is a closed
\emph{characterisation}, and it reduces the level to a search of $119{,}509$ nodes in place of the
original enumeration. That the two independent routes agree on $49{,}961$ is the point: the pins are
not merely necessary conditions that happen to hold, they cut out the level exactly, and nothing
outside them is doing any work. The evaluation also reports no subset sum within $10^{-12}$ of the
threshold, so the floating comparison is measured to be safe rather than assumed. For orientation,
the baseline $20$ smallest candidates carry $0.0907260944\ldots$, exceeding the threshold by
$T_{59}-2=0.0023501514\ldots$, and every admissible $S$ is a perturbation of that baseline within
this budget.

\begin{proposition}[How many rungs the ladder has]\label{prop:rungcount}
Let $(a,b)$ be a Pythagorean pair with $a<b$, $t=b/a$ and $|U|=n$. Positivity of the two masses gives
$-t<k\le0$, so $k$ takes at most $\lceil t\rceil$ values, and $t\le t_n$ where $t_n+1/t_n=T_n$. Hence
the number of rungs available at level $n$ is at most $\lceil t_n\rceil$, and
\[ t_n>m\iff T_n>m+\tfrac1m . \]
In particular the ladder has exactly the two rungs $k\in\{0,-1\}$ at every level $n\le1412$, and the
third rung $k=-2$ opens only at $n=1413$.
\end{proposition}

Two readings. First, the hypothesis $\sigma(ab)<5/2$ under which
Proposition~\ref{C-cor:kdetermined} confines $k$ to $\{0,-1\}$ is not a restriction anywhere a search
operates: it is equivalent to $t<2$, hence to $n\le1412$, and at level $59$ one has $t=1.049668$ with
$\lceil t\rceil=2$. The confinement is unconditional throughout the range the rest of this note
inhabits, and the hypothesis only becomes real past level $1412$.

Second, $1413$ is not a new constant. It is the level at which $T_n$ reaches $5/2$, and that is the
same threshold as the all--odd case of Theorem~\ref{K-thm:parity}: an all--odd cycle needs mass $2$ on
primes $\ge3$, hence $T_n\ge2+\tfrac12$. The level at which the ladder acquires a third rung and the
level at which an all--odd cycle becomes possible are the same level, for the same reason
\textup{(\texttt{code/level\_structure.rs})}. The fourth rung would need $T_n>10/3$, beyond
$\pi(4\times10^{7})$.

\begin{proposition}[The admissible supports are not a tree]\label{prop:nottree}
There are admissible supports at level $60$ with no admissible subset at level $59$. Explicitly,
\[ U\;=\;\{p:p\le271\}\ \cup\ \{809,\,811\} \]
has $|U|=60$ and $\sigma(U)=2.0012091\ldots>2$, while every $59$--element subset has mass below $2$.
\end{proposition}
\begin{proof}
Deleting the largest prime loses the least mass, so if that deletion is inadmissible every deletion
is. Here $\sigma(U)-1/811=1.9999761\ldots<2$.
\end{proof}

Such supports are exactly those of the form $V\cup\{q\}$ with $|V|=59$, $\sigma(V)\le2$,
$q>\max V$ and $q<1/(2-\sigma(V))$, and they are not exceptional. Taking $V=\{p\le271\}\cup\{r\}$
with $r>794$, so that $\sigma(V)<2$, each $r$ admits every prime $q$ in
$\bigl(r,\ 1/(2-\sigma(V))\bigr)$: at $r=797$ that interval reaches $190{,}414$ and contains
$17{,}061$ primes, and summing over $r<4000$ gives $60{,}831$ such supports from this one shape alone
\textup{(\texttt{code/level\_structure.rs})}.

The consequence is for the organisation of a level search, and it is independent of
Proposition~\ref{K-prop:levelbarrier}. That proposition says level $60$ has infinitely many admissible
supports; this one says that even the part reachable with bounded primes is not exhausted by the
base--plus--tail organisation, because the admissible supports do not form a tree under adjunction.
Closing every one--new--prime family at level $60$, which is question~(i) of
Section~\ref{K-sec:status} of the core note, would therefore still not close level $60$. The searches of
Propositions~\ref{A-prop:tailkill} and~\ref{K-prop:close59} are stated over one--new--prime families
throughout and are unaffected; what changes is the reading of what completing them would buy.

\begin{proposition}[The level--minimal structure at any $c$]\label{prop:levelstructure}
Let $U$ satisfy $\sum_{p\in U}1/p\ge c$ with $|U|=n(c)=\min\{n:T_n\ge c\}$. Then every prime $p$ with
$1/p\ge T_{n(c)+1}-c$ lies in $U$, and $\max U<1/\bigl(c-T_{n(c)-1}\bigr)$, the denominator being
positive by minimality of $n(c)$.
\end{proposition}
\begin{proof}
Both are Propositions~\ref{prop:window} and~\ref{prop:band} with $2$ replaced by $c$; neither
argument uses the value.
\end{proof}

The slack $T_{n(c)}-c$ is the total mass a support may waste, and it collapses as $c$ grows, so the
support becomes more forced, not less \textup{(\texttt{code/level\_structure.rs})}:
\[
\begin{array}{l|ccccc}
c & 2 & 9/4 & 7/3 & 5/2 & 3\\\hline
n(c) & 59 & 231 & 400 & 1{,}413 & 361{,}139\\
T_{n(c)}-c & 2.35\!\cdot\!10^{-3} & 2.33\!\cdot\!10^{-4} & 2.44\!\cdot\!10^{-4} &
  2.07\!\cdot\!10^{-5} & 3.09\!\cdot\!10^{-8}\\
\text{forced primes} & 39 & 181 & 259 & 1{,}175 & 314{,}407\\
\text{forced fraction} & 66\% & 78\% & 65\% & 83\% & 87\%\\
\max U< & 793.7 & 2{,}194 & 8{,}284 & 15{,}592 & 6.19\!\cdot\!10^{6}
\end{array}
\]
The row $c=2$ is Propositions~\ref{prop:window} and~\ref{prop:band}, recovered. The row $c=3$ is the
three--cycle of Proposition~\ref{K-prop:kcycles}: at its minimal level a three--cycle contains
\emph{every} prime up to $4{,}476{,}608$, that is $314{,}407$ of its $361{,}139$, and no prime beyond
$6.19\times10^{6}$. Its support is almost completely determined, with a slack of $3.09\times10^{-8}$
to distribute among the remaining $46{,}732$ places.

The same reading applies with the mass restricted to primes $\ge P$, which is the setting of
Proposition~\ref{K-prop:twoparam}. At $P=3$, the all--odd case of Theorem~\ref{K-thm:parity}, the
minimal level is $n=1412$ with slack $2.069\times10^{-5}$; every odd prime up to $9{,}484$ is forced,
$1{,}174$ of the $1{,}412$ or $83\%$, no prime exceeds $15{,}592$, and the product is $10^{5040.1}$.
At $P=5$ the figures are $n=40{,}259$, $24{,}770$ forced, $\max U<1.61\times10^{6}$ and $10^{209582}$,
reproducing Proposition~\ref{K-prop:twoparam} from the structural side.

That the all--odd case is $83\%$ forced does not make it a feasible search, and it is worth
recording why, since the shape invites the comparison with level $59$. There the free part was $20$
primes from $99$ candidates and the count was $49{,}961$. Here it is $238$ from $643$: the forced
mass is $1.97751561\ldots$, so the free primes must carry more than $0.02248438\ldots$ against a
baseline of $0.02250507\ldots$, an excess of $2.069\times10^{-5}$. A traversal capped at
$4\times10^{8}$ nodes reaches $1.977\times10^{8}$ admissible supports without exhausting the space
\textup{(\texttt{code/level\_structure.rs})}. The slack per free place, not the forced fraction, is
what decides enumerability, and here it is $87$ times larger per place than at level $59$. Question
(ii) of this section does not close this way.

\begin{proposition}[The barrier stratified by the smaller support]\label{prop:split}
Let $(a,b)$ be a two--cycle and put $k=\min(\omega(a),\omega(b))$. Then
\[ k=1\ \Rightarrow\ |U|\ge1413, \qquad k=2\ \Rightarrow\ |U|\ge69, \qquad
   k\ge3\ \Rightarrow\ |U|\ge59, \]
and the last is attained. Consequently every two--cycle with $|U|\le68$ has
$\min(\omega(a),\omega(b))\ge3$.
\end{proposition}
\begin{proof}
Write $T=T_{|U|}$. The supports are disjoint, so $\sigma(a)+\sigma(b)\le T$, and
$\sigma(a)\le T_k$ for the member with $k$ primes. The product $x(T-x)$ is increasing for
$x\le T/2$, so $\sigma(a)\sigma(b)\le m(T-m)$ with $m=\min(T_k,T/2)$, and $\sigma(a)\sigma(b)=1$
forces $m(T-m)\ge1$. For $k\ge3$ one has $T_k\ge T_3=31/30>T/2$ in the relevant range, so $m=T/2$
and the condition is $T\ge2$, i.e.\ $|U|\ge59$. For $k\le2$ one has $T_k<T/2$, so $m=T_k$ and the
condition is $T\ge T_k+1/T_k$: at $k=1$, $T\ge2.5$ and $|U|\ge1413$; at $k=2$, $T\ge2.0333\ldots$ and
$|U|\ge69$.
\end{proof}

The stratification is sharp at each end. At $k=1$ the member is a prime $p$, so $\sigma(b)=p\ge2$ and
$b$ must carry mass $2$ on primes other than $p$, which needs $1412$ of them; at $k=2$ the best case
is $a=6$, whence $\sigma(b)\ge6/5$ on primes $\ge5$, which needs $67$. This contains and extends
Proposition~16 of \cite{Kovic2012}, that no two--cycle has both members a product of two primes: the
present statement rules out \emph{either} member having two prime factors unless $|U|\ge69$, and
either having one unless $|U|\ge1413$. Since Proposition~\ref{K-prop:close59} gives $|U|\ge60$, the
window $60\le|U|\le68$ that any level--by--level search must cross admits only splits with both
$\omega(a)\ge3$ and $\omega(b)\ge3$ \textup{(\texttt{code/split\_bound.gp})}.

\begin{proposition}[The derivative cannot outgrow its own support]\label{prop:growth}
Let $a_0\to a_1\to\dots\to a_L$ be an orbit segment of the arithmetic derivative with every $a_i$
squarefree, put $G=a_L/a_0$ and let $U=\bigcup_{i<L}\operatorname{supp}(a_i)$. Then
\[ \sum_{p\in U}\frac1p\ \ge\ \frac{L}{\lceil L/2\rceil}\,G^{1/L},
\qquad\text{equivalently}\qquad
G\ \le\ \Bigl(\frac{\lceil L/2\rceil}{L}\,T_{|U|}\Bigr)^{\!L}, \]
and for even $L$ the geometric growth rate satisfies $G^{1/L}\le T_{|U|}/2$.
\end{proposition}
\begin{proof}
Squarefreeness gives $\gcd(a_i,a_i')=1$, so consecutive members are coprime and
$S_p=\{i<L:p\mid a_i\}$ has no two consecutive elements: it is independent in the path on $L$
vertices, whence $|S_p|\le\lceil L/2\rceil$. Then $\sum_{i<L}\sigma(a_i)=\sum_{p\in U}|S_p|/p\le
\lceil L/2\rceil\sum_{p\in U}1/p$, while $\prod_{i<L}\sigma(a_i)=a_L/a_0=G$ gives
$\sum_{i<L}\sigma(a_i)\ge LG^{1/L}$ by AM--GM. The second form inverts the first using
$\sum_{p\in U}1/p\le T_{|U|}$.
\end{proof}

This is the general shape of which the rest of the dictionary is a specialisation: at $L=2$ it reads
$\sum_{p\in U}1/p\ge2\sqrt G$, which is Corollary~\ref{cor:cost} with $G=r$; at $G=1$ it returns
$c=L/\lceil L/2\rceil\to2$, which is Proposition~\ref{K-prop:kuniform} on the path in place of the
cycle. The statement it adds is dynamical rather than Diophantine, and it bounds the growth branch
of Ufnarovski--\AA hlander's Conjecture~2: an orbit may tend to infinity, but only slowly unless it
accumulates an enormous support \textup{(\texttt{code/growth\_bound.gp})}:
\[
\begin{array}{l|cccccc}
\text{rate }g & 1.05 & 1.10 & 1.25 & 1.50 & 1.75 & 2\\\hline
|U|\ \ge & 97 & 170 & 1{,}413 & 361{,}139 & 4.6\times10^{9} & 4.3\times10^{16}\\
\text{largest prime}\ \ge & 5.4\!\cdot\!10^{2} & 1.0\!\cdot\!10^{3} & 1.2\!\cdot\!10^{4} &
 5.2\!\cdot\!10^{6} & 1.2\!\cdot\!10^{11} & 1.8\!\cdot\!10^{18}
\end{array}
\]
A squarefree orbit that doubles at every step needs some $4.3\times10^{16}$ distinct primes in its
supports and a prime past $1.8\times10^{18}$: sustained rapid growth is not merely unobserved but
impossible on a small support. The hypothesis that every $a_i$ is squarefree is essential and is not
generic along orbits; it is automatic on cycles, by Ufnarovski--\AA hlander, which is why the cycle
rows of the dictionary carry no such caveat.

\begin{corollary}[Sustained growth needs scale, and the scale is enormous]\label{cor:growthscale}
Let $a_0\to\dots\to a_L$ be a squarefree orbit segment with every member at most $X$ and geometric
growth rate at least $g>1$. Then
\[ \log_2X\ \ge\ \sqrt{\,n(2g)\,\log_2 g\,},\qquad n(c)=\min\{n:T_n\ge c\}. \]
\end{corollary}
\begin{proof}
Proposition~\ref{prop:growth} gives $\sum_{p\in U}1/p\ge2g$, so $|U|\ge n(2g)$. Every member is at
most $X$, so $\omega(a_i)\le\log_2X$ and $|U|\le\sum_{i<L}\omega(a_i)\le L\log_2X$; growth at rate
$g$ from $a_0\ge2$ gives $a_L\ge g^{L}$, so $L\le\log_gX=\log_2X/\log_2g$. Combining,
$(\log_2X)^2/\log_2g\ge n(2g)$.
\end{proof}

The threshold is explosive \textup{(\texttt{code/growth\_bound.gp})}:
\[
\begin{array}{l|ccccc}
g & 1.10 & 1.25 & 1.50 & 1.75 & 2\\\hline
\text{members must exceed} & 10^{1.5} & 10^{6.4} & 10^{138} & 10^{18420} &
 10^{6.23\times10^{7}}
\end{array}
\]
The arithmetic derivative can therefore grow slowly at any scale and quickly only at scales no
computation reaches. Rate $1.25$ needs members past $2.6\times10^{6}$, so it is visible inside the
sweep of Proposition~\ref{prop:nearmiss}; rate $1.5$ needs $10^{138}$; and a squarefree orbit cannot
double at every step unless its members exceed $10^{6.2\times10^{7}}$.

Inverting the corollary bounds the rate itself. The largest $g$ a squarefree orbit can sustain up to
$X$ solves $n(2g)\log_2g=(\log_2X)^2$, and since $n(c)$ grows like
$\operatorname{Li}\bigl(e^{e^{c-M}}\bigr)$ this gives $e^{e^{2g}}\asymp(\log_2X)^2$, that is
\[ g\ \sim\ \tfrac12\log\log\log X . \]
\[
\begin{array}{l|cccccc}
X & 10^{8} & 10^{20} & 10^{100} & 10^{1000} & 10^{10^{5}} & 10^{10^{7}}\\\hline
\text{max sustainable }g & 1.2768 & 1.3590 & 1.4808 & 1.6174 & 1.8134 & 1.9537
\end{array}
\]
The ceiling is a triple logarithm, so it is effectively a constant: even at $X=10^{10^{7}}$ the rate
stays below $2$, and for any range a computation can address it is below $1.5$. The reading for the
escape branch of Ufnarovski--\AA hlander's Conjecture~2 is that an orbit tending to infinity while
remaining squarefree must do so sub--geometrically in every bounded window, its rate decaying like
the triple logarithm of that window. Escape is possible; escape at a fixed geometric rate is not.

The mass window transfers to the normal form as well. In Proposition~\ref{K-prop:oneprime} a solution
is $a=Mp$, so $\sigma(a)=\sigma(M)+1/p$, and Proposition~\ref{prop:window} places $\sigma(a)$ in
$[1/t_n,t_n]$. Hence
\[ \sigma(M)\ \in\ \bigl[\,1/t_n-1/p,\ t_n\,\bigr], \]
which at level $59$ is $[0.9527-1/p,\ 1.0497]$: the base of the one--prime normal form carries mass
within five per cent of $1$, just as each member does, and $\omega(M)\ge2$ by
Proposition~\ref{prop:split}.

The region is not unreachable, only unsweepable. The witness $n_0$ of
Theorem~\ref{thm:lyapfalse} has $n_0$ and $n_0'$ both squarefree and
$r(n_0)=\sigma(n_0)\sigma(n_0')=1.03944\ldots$, so $n_0''>n_0$: the value $1$ that a solution must
hit is interior to the range of $r$, not extremal, and no inequality on $r$ alone can settle \#307.
The same construction approaches $1$ from below. Tuning $\sigma(a)$ against the forced
$\sigma(a')=\tfrac{31}{30}+\varepsilon$ gives
\[ a=3359\!\!\prod_{\substack{7\le p\le269\\p\notin\{229,271\}}}\!\!p,\qquad
   a'=2\cdot3\cdot5\cdot383\cdot113717\cdot P_{99}, \]
with $P_{99}$ prime, both $a$ and $a'$ squarefree, $108$ digits, and
\[ r(a)=0.99207796\ldots,\qquad |r(a)-1|=7.9220\times10^{-3}, \]
against the smallest relative defect $0.4465$ over the swept range
\textup{(\texttt{code/nearmiss\_tune.gp})}.

The same tuning with its acceptance window mirrored, $\sigma(a)$ just \emph{above} $30/31$ rather than
just below, reaches $1$ from the other side and closer. With
\[ a=3301\!\!\prod_{\substack{7\le p\le293\\ p\notin\{137,263\}}}\!\!p, \qquad
   a'=2\cdot3\cdot5\cdot P_{117}, \]
where $P_{117}$ is a $117$--digit prime, one has $a$ of $118$ digits and, exactly,
\[ r(a)=1.001014244\ldots,\qquad |r(a)-1|=1.0142\times10^{-3},\qquad a''>a. \]
This is exact rather than a lower bound, since $a'$ factors completely: the cofactor after trial
division is prime, so $\sigma(a')$ is known and not merely bounded below. It improves the record on
both sides at once, against $7.9220\times10^{-3}$ from below and $3.944\times10^{-2}$ from above,
the latter being $r(n_0)$ of Theorem~\ref{thm:lyapfalse}
\textup{(\texttt{code/nearmiss\_above.gp}, re--checked independently in
\texttt{code/nearmiss\_above\_verify.gp})}.

That record is not the natural limit, and seeing why identifies what the near--miss problem actually
is. In every witness of this shape $a'=2\cdot3\cdot5\cdot P$ with $P$ prime, so
$\sigma(a')=\tfrac{31}{30}+P^{-1}$ with $P^{-1}$ of size $10^{-118}$, and therefore
\[ |r(a)-1|=\tfrac{31}{30}\,\bigl|\sigma(a)-\tfrac{30}{31}\bigr| \]
to any precision that matters. The quantity $a'$ has dropped out. What remains is the Diophantine
question of how closely a sum of \emph{distinct prime reciprocals} approaches $30/31$ from above, and
the $10^{-3}$ above is only the granularity of adjusting by a single prime below $4000$.

Closing the deficit in stages removes that ceiling. With $d=\tfrac{30}{31}-\sigma(S)$ for
$S$ the primes in $[7,271]$, take $p_1$ the least prime exceeding $1/d$, then $p_2$ the least
exceeding $1/(d-p_1^{-1})$, and finally $p_3$ the greatest prime below the reciprocal of what
remains; the overshoot is then of order $\mathrm{gap}(p_3)/p_3^{2}$. Each rung squares the precision.
Sweeping the choice at the first two rungs and retaining those that satisfy the forcing congruence
and have $a'=2\cdot3\cdot5\cdot P$ gives, with
$(p_1,p_2,p_3)=(569,1949,15461)$,
\[ a=569\cdot1949\cdot15461\!\!\prod_{7\le p\le271}\!\!p \quad(58\text{ primes},\ 120\text{ digits}),
   \qquad a'=2\cdot3\cdot5\cdot P_{118}, \]
and, exactly,
\[ r(a)=1+6.278713\times10^{-9},\qquad a''>a. \]
This is $1.6\times10^{5}$ times closer to $1$ than the preceding record.

That record is on the wrong side of $1$, and identifying the family says why. Every witness here has
$a'=2\cdot3\cdot5\cdot P$ with $P$ prime, which is exactly the stratum $M=30$ of
Proposition~\textup{\ref{K-prop:oneprime}}: with $M'=31$ the equation $MN-M'N'=M^2$ reads
$30N-31N'=900$, and the proposition's $N=M+M'p$, $N'=Mp$ become $a=31P+30$, $a'=30P$. Two identities
follow, both exact: $a''=(30P)'=31P+30$, so $r=a''/a=(31P+30)/a$, and a solution of the stratum is
precisely $30a-31a'=900$. Hence
\[ \sigma(a)=\frac{30P}{31P+30}=\frac{30}{31}-\frac{900}{31a}\;<\;\frac{30}{31} \]
for every solution of this stratum, strictly. Approaches to $1$ from \emph{above}, which is what the
$6.28\times10^{-9}$ above and the $3.76\times10^{-9}$ of a wider sweep are, therefore lie on a side
the stratum cannot populate. They stand as records for the spectrum of $r$ and cannot converge on a
solution.

Aimed below instead, where the stratum's solutions must lie, the same ladder gives
\[ a=431\cdot68111\cdot2792213\!\!\prod_{7\le p\le271}\!\!p \quad(58\text{ primes},\ 123\text{ digits}),
   \qquad a'=2\cdot3\cdot5\cdot P_{122}, \]
and, exactly,
\[ r(a)=1-1.654663\times10^{-13},\qquad \sigma(a)<\tfrac{30}{31},\qquad a''<a, \]
which is $4.8\times10^{10}$ times closer to $1$ than the previous record from below. For this witness
$30a-31a'$ is a $112$--digit number where a solution needs $900$, which is the honest measure of the
distance remaining: the ratio is close, the integer is not.

Neither construction is exhausted. Over the same sweeps the smallest \emph{available} deficit was
$2.043\times10^{-18}$ from below and $1.209\times10^{-14}$ from above; what limits the reported values
is the rate at which $a'$ takes the form $2\cdot3\cdot5\cdot P$, not the tuning
\textup{(\texttt{code/nearmiss\_ladder.gp}, \texttt{code/nearmiss\_below.gp}, re--checked in
\texttt{code/nearmiss\_ladder\_verify.gp} and \texttt{code/nearmiss\_below\_verify.gp})}.

We stop here, and the reason is a measurement rather than a budget. Between the two from--below
witnesses the ratio improved by three orders of magnitude while $30a-31a'$ fell from $114$ digits to
$112$, against the $3$ a solution requires. Exhausting the available deficit would gain perhaps five
more digits of the $112$. The near--miss frontier measures $|r-1|$, a ratio; membership of the
stratum is the integer identity $30a-31a'=900$, and the two are not commensurable at this scale.
Driving the ratio further is therefore not evidence about the integer, and the line is closed as a
line of attack. What survives it is structural and is stated above: the family is the $M=30$ stratum
of Proposition~\ref{K-prop:oneprime}, its solutions satisfy $\sigma(a)<30/31$ strictly, and approaches
from above cannot reach one. Those three facts are machine--checked in
\texttt{lean/Erdos307/StratumM30.lean}.

The constant $30/31$ is not special. In the normal form of Proposition~\ref{K-prop:oneprime} a solution
of the $M$--stratum is $a=Mp$, $b=N=M+M'p$, so $\sigma(a)=\sigma(M)+1/p$ and, by the identity
$(\sigma(M)+1/p)\sigma(N)=1$ of part~(b), $\sigma(b)=1/\sigma(a)<1/\sigma(M)=M/M'$, with defect
exactly $M^2/(M'(M+M'p))$. Hence \emph{every} stratum is one--sided: for no squarefree $M$ does the
half--space $\sigma(b)\ge M/M'$ contain a solution, and $M=30$ recovers $30/31$. This is a corollary
of Proposition~\ref{K-prop:oneprime}(b) rather than a new result, but it is the constraint that governs
search direction, so it is machine--checked in \texttt{lean/Erdos307/StratumGeneral.lean}. The
consequence for the near--miss programme is that the mirrored tuning of the previous paragraph is not
merely unlucky at $M=30$; the corresponding tuning is empty in every stratum. Consistently with the two--sided statement above, $a$
lies past the barrier: $118>113$ digits, and $r$ sits inside the barred neighbourhood
$(0.998740,\infty)$.

What remains unmeasurable is the density of $r$ at $1$,
since constructed points do not estimate a density, and that is what the heuristic constant needs.
The frontier reported above is therefore superseded as a record and unchanged as an argument.

The situation is not the usual one. Heuristics of Bateman--Horn or Hardy--Littlewood
type are normally checked against data, and the agreement is what justifies confidence in them. Here
the prediction that solutions exist rests on a constant that cannot be estimated by sampling,
because the barrier and the sampleable range are disjoint. The expected-count heuristic for \#307 is
therefore unverifiable in the usual way.

It is not, however, beyond computation, and the reason is that the constant is the barrier.

\begin{proposition}[The heuristic constant is the barrier's own probability]\label{prop:f1}
Model a random squarefree $a$ by $\mathbb{P}(p\mid a)=1/(p+1)$, and $a'$ by the coprimality model
that Proposition~\textup{\ref{prop:anticorr}} validates, so that $\mathbb{P}(q\mid a')=1/q$ for
$q\nmid a$. Then each prime independently lies in $a$, in $a'$, or in neither, with probabilities
$\tfrac1{p+1},\tfrac1{p+1},1-\tfrac2{p+1}$, and
\[ \sigma(a)+\sigma(a')=\sum_{p\,\mid\,aa'}\frac1p,\qquad \mathbb{P}(p\mid aa')=\frac2{p+1}. \]
Since a two-cycle forces $\sigma(a)\sigma(a')=1$ and hence $\sigma(a)+\sigma(a')\ge2$ by
\textup{AM--GM}, the density $f$ of $r=\sigma(a)\sigma(a')$ is supported in the event
$B:\sigma(a)+\sigma(a')\ge2$, so $f(1)\le C\,\mathbb{P}(B)$ with $C$ the conditional density of $r$
at $1$ given $B$. Numerically
\[ \mathbb{P}(B)=2.74\times10^{-90}. \]
\end{proposition}

The point is the order and not the digits. $\mathbb{P}(B)$ is a \emph{tail} probability of a
one-dimensional sum, not a density, which is what makes it stable: it moves from $4.09$ to $3.01$ to
$2.74\times10^{-90}$ as the grid is refined from $2\times10^4$ to $8\times10^4$ bins and the prime
range from $10^4$ to $10^6$, while a direct attempt at $f(1)$ on a two-dimensional grid moves by four
orders of magnitude per refinement and converges to nothing. Two checks support it. The same
computation with $\mathbb{P}(p\mid n)=1/p$ returns $\mathbb{P}(\sigma(n)\ge1)=0.0420$, reproducing
the measured density of prime-abundant integers to four places; and the single configuration in which
$aa'$ absorbs exactly the first $59$ primes has probability $\prod_{p\le277}2/(p+1)=10^{-96.0}$,
below $\mathbb{P}(B)$ by the factor one expects from the many other prime sets of mass $2$
\textup{(\texttt{code/barrier\_probability.rs})}.

One more thing follows, and it disarms the near-miss spectrum entirely.

\begin{proposition}[The near-miss infimum is zero, so closeness is not evidence]\label{prop:infzero}
\[ \inf\Bigl\{\,\bigl|\sigma(P)\sigma(Q)-1\bigr|\ :\ P,Q\ \text{finite disjoint prime sets}\,\Bigr\}=0, \]
and \#307 asks exactly whether the infimum is attained.
\end{proposition}
\begin{proof}
Split the primes into two infinite classes, say by parity of index. Each class has terms tending to
$0$ and divergent reciprocal sum, so by the classical subseries theorem each carries a subseries
summing to exactly $1$. This produces disjoint \emph{infinite} $P_\infty,Q_\infty$ with
$\sigma(P_\infty)=\sigma(Q_\infty)=1$. Truncating both at height $y$ gives finite disjoint sets with
$\sigma(P_y)\sigma(Q_y)\to1$.
\end{proof}

The analytic content is formalised, and in a slightly stronger form than the proof above uses. For
the infimum one does not need the subseries theorem, only that each mass can be placed in
$(1-\delta,1)$: \texttt{Erdos307.exists\_block\_sum\_near} takes the first consecutive block whose
sum exceeds $T-\varepsilon$ and observes that it overshoots by less than one term, and
\texttt{infimum\_zero} assembles two such blocks into a product within $\varepsilon$ of $1$. Both
carry $0$ \texttt{sorry} and the three standard axioms
\textup{(\texttt{lean/Erdos307/InfZero.lean})}. The arithmetic input, that the primes split into
two infinite classes each with reciprocal terms tending to $0$ and unbounded partial sums, is cited
rather than reproved, as in Theorem~\ref{thm:halflyap}.

The greedy form of that argument reaches $|\sigma(P)\sigma(Q)-1|=6.50\times10^{-9}$ with $144$
primes, against $0.4465$ for the swept frontier of Proposition~\ref{prop:nearmiss}
\textup{(\texttt{code/infimum\_zero.gp})}. So the analytic half of \#307 is empty: the mass equation
can be approximated as closely as one likes, and no near-miss, however good, is evidence that a
solution exists.

What the near-misses do measure is a trade-off, and it is the honest picture of the obstruction. A
solution must kill two defects at once, the mass defect $|\sigma(P)\sigma(Q)-1|$ and the size defect
$\log_{10}\bigl(\prod Q\big/\sigma(P)\prod P\bigr)$, the latter vanishing exactly when
$\prod Q=N_P$, which is Theorem~\ref{K-thm:structure}. Each is reachable alone and neither
construction moves the other:
\[
\begin{array}{lcc}
\text{construction} & \text{mass defect} & \text{size defect}\\[2pt]
\text{greedy, } Q \text{ free} & 6.50\times10^{-9} & 10^{341.3}\\
Q=\text{primes}(a'),\ \text{\texttt{nearmiss\_tune.gp}} & 7.92\times10^{-3} & 0
\end{array}
\]
Fixing $Q$ as the prime set of $a'$ makes the size defect vanish identically and leaves the mass
defect at $10^{-3}$; choosing $Q$ freely drives the mass defect to $10^{-9}$ and costs $341$ orders
of magnitude in size. \#307 is the demand that both columns read zero simultaneously, and nothing in
this paper moves them together.

\begin{remark}[What the calibrated heuristic says, and it is not encouraging]\label{rem:calibrated}
With $\mathbb{E}(X)\sim f(1)\,\delta\log X$ and $\delta=0.3721$, Proposition~\ref{prop:f1} gives an
expected number of two-cycles below the proved barrier of about $10^{-88}$, and $\mathbb{E}(X)$ does
not reach $1$ until $X$ is of order $10^{4\times10^{89}}$. So the heuristic favours \textsc{yes} only
in the limiting sense that $\log X$ diverges; at every scale that any search or any barrier argument
will reach, it predicts nothing. Two consequences, both negative and both worth stating. A proof of
the negative direction would have to contradict a divergence, and a search for a solution is
hopeless not merely at present but at any scale expressible in the notation of this paper. The
constant is small for a reason that is not incidental: it \emph{is} the barrier, seen as a
probability rather than as a counting bound, which is the third independent appearance of the same
Mertens threshold in this work.

The label is \textup{HEURISTIC} and the model is doing the work. What is claimed is only that the
constant, said above to be unavailable to sampling, is available to computation once the coprimality
model of Proposition~\ref{prop:anticorr} is granted.
\end{remark}

\begin{remark}[What is left, on both sides]\label{rem:whatisleft}
The two directions can now be specified rather than merely called open.

\emph{Negative.} Theorem~\ref{thm:noinvariant} closes the certificate route on both branches, so a
proof must count rather than certify; and by Remark~\ref{rem:whatcounting} the count it must defeat
diverges. A proof of the negative direction would therefore have to contradict a divergent expected
count, which is to say that the direction is not merely unproved but probably false.

\emph{Positive.} Three routes are available in principle, and two are closed quantitatively.
\begin{enumerate}[label=\textup{(\alph*)},leftmargin=2.2em]
\item \emph{Exhibit a witness.} Excluded by Remark~\ref{rem:calibrated}: the first is expected near
$10^{4\times10^{89}}$.
\item \emph{Bound the count below by analytic means.} Excluded by size. At the proved barrier the
main term $f(1)\delta\log X$ lies $10^{200.6}$ below the trivial count $X$, and the saving a method
would need in order to keep the main term above its own error grows like $X$ itself. Sieve and
circle-method errors save a power of $X$ or a power of $\log X$. So an analytic proof here would
have to be \emph{exact}, with error identically zero, rather than sharp.
\item \emph{Construct exactly.} This is the only branch alive, and it exists: the Pell orbit of
Remark~\ref{A-rem:lehmer}. It is alive but narrower than it was, and the narrowing should not be
mistaken for proximity. Two cautions, both quantitative. The searches of
Remark~\ref{A-rem:notlehmer} came back empty over $3.99\times10^{11}$ candidates, but empty was the
\emph{predicted} outcome: the expected count was $2.4\times10^{-56}$, so the computation confirms the
heuristic model rather than deciding anything, and a confirmed model is not a decided problem. And by
Proposition~\ref{K-prop:levelbarrier} those searches cannot be continued upward at all, since the
admissible supports are infinite at every level past $59$. The unconditional exclusions of this note
stop where enumeration stops, which is now.
Proposition~\ref{K-prop:tailbound} bounds the tail prime by $4N$, so the orbit contributes nothing past
a bounded initial segment and the construction is a finite search over the divisors of $4D$ rather
than a Lehmer primality question; Remark~\ref{A-rem:notlehmer} then makes that finite search
heuristically empty on all $34$ immune families, by two independent estimates. What survives is the
branch itself, not its expectation: the construction must now either come from a base outside the
immune list or exploit the parametrisation $q=2(N\pm k)/m^2$ directly.
\end{enumerate}

The two exclusions interlock, and the interlock is the real statement. Exactness is what algebra
supplies, and Theorem~\ref{thm:ff} excludes every argument that survives the passage to $K[t]$,
where the conclusion is false. The Pell construction escapes that exclusion for one reason only:
its final input is the primality of a determined term, and primality does not survive to $K[t]$. So a
successful method must be \emph{archimedean and exact at the same time}, a combination which is
unusual but not unprecedented, continued fractions solving Pell being the standard instance, and
Remark~\ref{A-rem:lehmer} is exactly the analogue here.

The labels differ across the three branches and the difference matters. The exclusions of (a) and
(b) are \textup{HEURISTIC}: both rest on the value of $f(1)$ from Proposition~\ref{prop:f1}, hence on
the coprimality model, and a reader who rejects that model recovers both branches. The interlock
itself is not heuristic: Theorem~\ref{thm:ff} and Proposition~\ref{K-prop:nopoly} are
\textup{PROVED}, and so is the reduction of branch (c) to a primality question. So what follows is
stated at the strength of its weakest link, which is the model.

What remains of \#307 is therefore one question and not a programme: whether some base admits a
Pythagorean tail prime, together with the sign selection $k=0$, which Remark~\ref{A-rem:lehmer} shows
is a choice between $\{0,-1\}$ rather than a coincidence of measure zero. By
Proposition~\ref{K-prop:nopoly} the orbit parametrising the Pythagorean locus is forced to be
exponential rather than polynomial. That was previously read as placing the question in the Lehmer
class rather than the Schinzel class, hence open for every non-degenerate sequence; but by
Proposition~\ref{K-prop:tailbound} a prime tail satisfies $q\le4N$, so the exponential orbit supplies
no candidate past a bounded initial segment and the question is neither Lehmer nor Schinzel. It is a
finite search over the divisors $m\mid4D$ with $AB+Dm^2$ a square, which
Remark~\ref{A-rem:notlehmer} makes heuristically empty on all $34$ immune families. The residue is
thus smaller and better located than it was, and still undecided: the bound is proved, the
emptiness is not, and no base outside the immune list has been examined at all. That is why nothing
in this paper decides the problem.
\end{remark}

The measurement does expose a sharp structural fact, which turns out to be exactly the barrier seen
statistically.

\begin{proposition}[The barrier is statistically complete]\label{prop:anticorr}
$\mathbb{E}[\sigma(a')]$ falls by an order of magnitude as $\sigma(a)$ grows, from $0.2035$
unconditionally to $0.0207$ on $\sigma(a)\ge1.3$. The whole of this suppression is coprimality.
Writing the prediction for an integer chosen at random subject to $\gcd(n,a)=1$, namely
$\sum_{p\nmid a}p^{-2}$, the ratio of measured to predicted is
\[ 0.65,\quad 0.82,\quad 0.91,\quad 0.92,\quad 0.95,\quad 1.01 \]
at $\sigma(a)\ge0,\,0.5,\,0.9,\,1.0,\,1.2,\,1.3$: it converges to $1$ precisely in the regime a
solution would inhabit \textup{(\texttt{code/a10\_anticorr.rs})}.
\end{proposition}

The mean is not the moment that $f(1)$ depends on. Since $f(1)$ is an integral of the joint density
of $(\sigma(a),\sigma(a'))$ along $\sigma(a)\sigma(a')=1$, a model could match the conditional mean
and still misstate the density; the next moment is therefore the sharper test. The model makes a
prediction for it too, and a per--$a$ one: given $\operatorname{supp}(a)$, each $q\nmid a$ lands in
$a'$ independently with probability $1/q$, so
\[ \mathbb{E}[\sigma(a')]=\sum_{q\nmid a}\frac1{q^{2}},\qquad
   \operatorname{Var}[\sigma(a')]=\sum_{q\nmid a}\frac1{q^{3}}\Bigl(1-\frac1q\Bigr). \]
Measured against these over all $a\le10^{7}$ with $a,a'$ squarefree
\textup{(\texttt{code/f1\_variance.rs})}, the ratio of measured to predicted variance is
\[ 1.111,\ 0.968,\ 0.894,\ 0.956,\ 0.938,\ 0.860,\ 0.948,\ 1.030 \]
at $\sigma(a)\ge0,\,0.5,\,0.7,\,0.9,\,1.0,\,1.1,\,1.2,\,1.3$, against $0.648,\dots,1.010$ for the
mean. The variance is predicted as well as the mean, it shows no systematic drift, and the two
converge together. The top band rests on $438$ values and is correspondingly noisy, so we claim only
the absence of a trend, not the precise value. What the test was designed to detect (a model
matching the first moment while its tail degrades, which would inflate $f(1)$ and could in principle
send the true value to $0$) does not occur.

So on the sampled range $a'$ is mass-poor for one reason only, that it is coprime to $a$ and $a$ has
absorbed the small primes, and there is no further effect to exploit. This is an independent
confirmation of Proposition~\ref{C-prop:nogain}: the barrier is not merely near-optimal as a bound, it
already contains the entire statistical structure of the obstruction, and no refinement by
correlation is available.

The range matters, and the statement is about averages. Every $a\le4\times10^8$ with $\sigma(a)>1$
is even, so the measurement never sees an $a$ that omits the small primes. Such $a$ exist, since the
prime reciprocals diverge, and for them $a'$ is free to absorb $2$, $3$ and $5$ and is not mass-poor
at all: Theorem~\ref{thm:lyapfalse} exhibits one with $\sigma(a')=\tfrac{31}{30}$. Proposition~\ref{prop:anticorr}
bounds an expectation on a range, and nothing here bounds a supremum off it.

\subsection*{A Lyapunov candidate for the negative direction}

Theorem~\ref{thm:ff} rules out cycles over $\mathbb{F}_q[t]$ by exhibiting a function that provably
decreases, namely the degree. Proposition~\ref{C-prop:noplace} shows no \emph{absolute value} on
$\mathbb{Q}$ plays that role. It does not exclude functions that are not absolute values, and the
following is the outcome of a search among such functions. It is a conjecture, not a theorem, and it
is recorded because it is falsifiable and because the mechanism protecting it is one of the paper's
own results.

\begin{definition}\label{def:lyap}
For $\lambda>0$ and squarefree $n$ set $\delta_\lambda(n)=\log n+\lambda\,\sigma(n)$.
\end{definition}

\begin{proposition}[A decreasing $\delta_\lambda$ would settle \#307 negatively]\label{prop:lyapworks}
If for some $\lambda>0$ one has $\delta_\lambda(n')<\delta_\lambda(n)$ for every squarefree $n$ with
$n'$ squarefree, then the arithmetic derivative has no cycle of any length among such $n$, and in
particular \#307 has answer \textsc{no}.
\end{proposition}
\begin{proof}
A cycle would give $\delta_\lambda(n)<\delta_\lambda(n)$.
\end{proof}

Since $n'=n\sigma(n)$, the requirement is the pointwise inequality
\[ \log\sigma(n)\;<\;\lambda\bigl(\sigma(n)-\sigma(n')\bigr). \tag{$\ast$} \]
Around a two-cycle the two instances of $(\ast)$ have increments summing to
$\log(\sigma(a)\sigma(b))=0$, so they cannot both hold: the admissible range of $\lambda$ closes
exactly at a cycle. The content is therefore whether that range is non-empty.

Below the barrier the criterion is not conjectural: it holds with the explicit value
$\lambda=\tfrac12$, by an elementary inequality.

\begin{proposition}[The obstruction is to the shape of $\delta_\lambda$, not to $\lambda$]\label{prop:lyapshape}
The closing of the admissible range of $\lambda$ is not special to a linear dependence on $\sigma$.
Let $F(n)=\log n+G(\sigma(n))$ for an arbitrary $G$. Since $n'=n\,\sigma(n)$, one step gives
\[ F(n')-F(n)\;=\;\log\sigma(n)+G\bigl(\sigma(n')\bigr)-G\bigl(\sigma(n)\bigr), \]
so around a two--cycle $(a,b)$ the two $G$--terms telescope and the increments sum to
$\log\bigl(\sigma(a)\sigma(b)\bigr)=0$ for \emph{every} $G$. A strictly decreasing $F$ of this shape
would settle \#307, and none can be exhibited without already excluding cycles: the $G$--part
cancels identically and only the mass identity survives. The relevant identity is
$\sigma(n)\sigma(n')=n''/n$, so $\sigma(n)\sigma(n')=1$ holds exactly when $n''=n$: the configuration
that defeats the class occurs at a cycle and nowhere else. Any Lyapunov argument must therefore use a
potential that sees more of $n$ than $\sigma(n)$ alone.
\end{proposition}

\begin{theorem}[The half-mass Lyapunov function]\label{thm:halflyap}
Put $\delta(n)=\log n+\tfrac12\sigma(n)$. Let $n$ be squarefree with $n'$ squarefree. If
$\sigma(n\,n')<2$ then $\delta(n')<\delta(n)$.
\end{theorem}
\begin{proof}
By Theorem~\ref{K-thm:structure} $n$ and $n'$ are coprime, and both are squarefree, so
$\sigma(n\,n')=\sigma(n)+\sigma(n')$. The hypothesis therefore gives $\sigma(n')<2-\sigma(n)$, hence
\[ \tfrac12\bigl(\sigma(n)-\sigma(n')\bigr)\;>\;\sigma(n)-1\;\ge\;\log\sigma(n), \]
the last step by $\log x\le x-1$. If $\sigma(n)=1$ the two ends read $0$ and the first inequality is
still strict, since $\sigma(n)+\sigma(n')<2$ forces $\sigma(n')<1$. Now $n'=n\sigma(n)$ gives
$\delta(n')-\delta(n)=\log\sigma(n)-\tfrac12(\sigma(n)-\sigma(n'))<0$.
\end{proof}
Verified over all $7{,}480{,}605$ squarefree $n\le2\times10^7$ with $n'$ squarefree lying in the
region: $0$ violations, least slack $0.242$ \textup{(\texttt{code/invent\_lyapunov.rs})}. The analytic content is formalised:
\texttt{Erdos307.log\_lt\_half\_sub} proves $\log s<(s-t)/2$ for $s>0$ and $s+t<2$,
\texttt{delta\_decreasing} gives the increment form, and
\texttt{two\_le\_add\_of\_mul\_eq\_one} is the \textup{AM--GM} step of
Corollary~\ref{cor:lyapbarrier}; all three carry $0$ \texttt{sorry} and only the three standard
axioms \textup{(\texttt{lean/Erdos307/HalfLyap.lean})}. The arithmetic inputs, that $n$ and $n'$ are
coprime and that $n'=n\sigma(n)$, are Theorem~\ref{K-thm:structure} and the definition, so the
statement is \textup{VERIFIED} modulo those.

\begin{corollary}[A Lyapunov proof of the barrier]\label{cor:lyapbarrier}
The arithmetic derivative has no cycle, of any length, all of whose consecutive pairs satisfy
$\sigma(n\,n')<2$. In particular a two-cycle $a'=b$, $b'=a$ with $a\neq b$ has
$\sigma(ab)=\sigma(a)+\sigma(b)\ge2\sqrt{\sigma(a)\sigma(b)}=2$, with equality only if
$\sigma(a)=\sigma(b)=1$ and hence $a=b$; so $\sigma(ab)>2$ and therefore $ab\ge\Pi_{59}$.
\end{corollary}
\begin{proof}
A cycle would give $\delta(n)<\delta(n)$ by Theorem~\ref{thm:halflyap}. The second statement is
AM--GM together with the mass estimate of Theorem~\ref{K-thm:barrier}.
\end{proof}

There are exactly two ways to exploit the fact that increments cancel around a cycle, and the
second is foreclosed. Around a two-cycle the two increments of any function sum to zero, so an
obstruction must pin down either their \emph{sign}, which is the Lyapunov route above, or their
\emph{value in a finite group}: if some $\varphi$ satisfied
$\varphi(n')-\varphi(n)\equiv c\pmod m$ for all $n$ with $2c\not\equiv0$, a two-cycle would give
$0\equiv2c$. No such $\varphi$ exists, and not for want of searching:
Proposition~\ref{A-prop:localcomplete} proves the cycle system soluble modulo every $m$, so any
obstruction expressible as a congruence is empty a priori. Measured, the increments of
$\omega(n)$, of $\#\{p\mid n:p\equiv3\ (4)\}$, of $n$, and of $n+\omega(n)$ each hit every residue
class modulo $2$, $3$ and $4$ \textup{(\texttt{code/invent\_lyapunov.rs})}. The sign route is
therefore the only one of the two available, which is another form of the thesis that the
obstruction is ordered--archimedean rather than congruential (Remark~\ref{rem:ordered}).

The sign route is available but not weaker than the problem. Both branches of the
dichotomy can now be closed at once.

\begin{theorem}[Both branches are closed]\label{thm:noinvariant}
Let $S$ be the set of squarefree $n\ge2$ with $n'$ squarefree.
\begin{enumerate}[label=\textup{(\arabic*)},leftmargin=2.2em]
\item \textup{(Finite-group branch.)} \textsc{Conjecture}, downgraded. The claim that there is no
$\varphi$ on $S$ and modulus $m$ with $\varphi(n')-\varphi(n)\equiv c\pmod m$ for all $n$ and
$2c\not\equiv0\pmod m$ is \emph{not} established. Proposition~\ref{A-prop:localcomplete} settles only
the moduli it covers, $m\le210$, and the general statement is equivalent to the target rather than
independent of it: on a functional graph with no periodic orbit one may build a constant-increment
$\varphi$ for any $(m,c)$ by assigning $\varphi$ along depth, so nonexistence for all $m$ is
equivalent to the existence of a periodic orbit. Branch (1) therefore has exactly the status of
branch (2), which Proposition~\ref{prop:lyapuniversal} already proves for the sign side. What is
proved is the mechanism recorded after the proof.
\item \textup{(Sign branch, not weaker than the target.)} A $\delta$ on $S$ with
$\delta(n')<\delta(n)$ throughout exists if and only if $n\mapsto n'$ has no periodic orbit in $S$;
that condition implies a negative answer to \#307 and is not implied by it.
\end{enumerate}
\end{theorem}
\begin{proof}
(1) is not proved; see the statement. Proposition~\ref{A-prop:localcomplete} makes the cycle system
soluble modulo every $m\le210$, which does not reach the quantifier the branch needs.
(2) is Proposition~\ref{prop:lyapuniversal}; a two-cycle is a periodic orbit, and periodic orbits of
length other than two are not excluded by a negative answer to \#307.
\end{proof}

The mechanism of branch (1) is formalised. \texttt{Erdos307.two\_nsmul\_eq\_zero\_of\_period\_two}
proves that a constant increment $c$ in any abelian group forces $2c=0$ at a point of period two,
\texttt{const\_increment\_iterate} gives the general $\varphi(D^{[k]}n)=\varphi(n)+kc$,
\texttt{no\_period\_two\_of\_two\_nsmul\_ne\_zero} is the certificate such a $\varphi$ would supply,
and \texttt{zmod\_two\_mul\_eq\_zero\_of\_period\_two} states it over $\mathbb{Z}/m$; all with $0$
\texttt{sorry} and on \texttt{propext} alone
\textup{(\texttt{lean/Erdos307/Congruential.lean})}. What is cited rather than formalised is the
\emph{emptiness} of the class, that no such $\varphi$ with $2c\neq0$ exists for the arithmetic
derivative, which is Proposition~\ref{A-prop:localcomplete} and is arithmetic. Branch (1) is therefore
\textup{VERIFIED} in its mechanism and \textup{PROVED} in its emptiness, and the two should not be
conflated.

The taxonomy that makes this a closure rather than two isolated facts is the displayed thesis above,
that an invariant whose increments cancel around a cycle can be exploited only through their sign or
their value in a finite group. That thesis is not a theorem, and an argument outside it is not
excluded by anything here.

This is a statement about methods, not about \#307, and it is the natural endpoint of the structural
programme: Proposition~\ref{C-prop:noplace} removes the absolute values, Theorem~\ref{thm:ff} removes
the arguments that survive to $K[t]$, Proposition~\ref{A-prop:localcomplete} removes the congruences,
and Theorem~\ref{thm:noinvariant} removes what is left of the pointwise route. A proof of the
negative direction, if there is one, must be global: it must count, and not certify.

\begin{remark}[What a counting proof would have to defeat]\label{rem:whatcounting}
The count it would have to defeat diverges. A two-cycle is exactly $r(a)=1$ for
$r(a)=\sigma(a)\sigma(a')=a''/a$, so if $r$ has a limiting density $f$ on the class of squarefree $a$
with $a'$ squarefree, then $\mathbb{P}(a''=a)=\mathbb{P}\bigl(r\in[1,1+\tfrac1a)\bigr)\sim f(1)/a$ and
\[ \mathbb{E}\bigl(\#\{a\le X:a''=a\}\bigr)\ \sim\ f(1)\,\delta\log X,\qquad \delta=0.3721\ldots, \]
$\delta$ being the measured density of the admissible class \textup{(\texttt{code/region\_shape.rs},
$3{,}721{,}148$ admissible $a\le10^7$)}. This diverges, so the heuristic predicts infinitely many
two-cycles, and any proof of the negative direction must account for a divergent expected count
producing no solutions at all.

The same formula locates them, and shows why the location is not knowable. Setting the expectation to
$1$ gives a first solution near $\exp\bigl(1/(\delta f(1))\bigr)$, which is exponential in $1/f(1)$:
at $f(1)=10^{-2}$ it is $10^{116.7}$, at $f(1)=5\times10^{-3}$ it is $10^{233.4}$, at $f(1)=10^{-3}$
it is $10^{1167}$. Halving the constant squares the answer. Consistency with
Theorem~\ref{K-thm:barrier} needs $f(1)<0.01034$, which is evidence of the softest kind, since an
expectation near $1$ with no solution observed is an ordinary event and not a contradiction. What
Proposition~\ref{prop:nearmiss} shows is that $f(1)$ is the one quantity here that cannot be sampled
\textup{(\texttt{code/heuristic\_location.gp})}. That $f(1)>0$ is not established either; the
constructed witnesses of Theorem~\ref{thm:lyapfalse} and of the near-miss frontier show only that $r$
takes values on both sides of $1$ and within $8\times10^{-3}$ of it.
\end{remark}

Corollary~\ref{cor:lyapbarrier} recovers the barrier by a route independent of the minimisation
used in Theorem~\ref{K-thm:barrier}: no extremality argument, no enumeration, only $\log x\le x-1$ and
the coprimality of Theorem~\ref{K-thm:structure}. It is also stronger in one direction, since it
excludes cycles of every length rather than two-cycles alone. What it does not do is decide
\#307: the region it covers is exactly the region where the barrier already says nothing can
happen, and above $\Pi_{59}$ the hypothesis $\sigma(n\,n')<2$ fails, which is where any solution
must live. The general conjecture is what remains.

Over all $153{,}397{,}754$ such $n\le4\times10^8$, $(\ast)$ holds for every
$\lambda\in(0.275438,\,1.250828)$ and fails for no step
\textup{(\texttt{code/lyap\_battery.rs})}. The interval narrows with the range, of width
$1.654$, $1.267$, $1.237$, $1.116$, $1.092$, $1.004$, $0.975$ at cutoffs $10^4$ through
$4\times10^8$;
the binding lower constraint is at $n=\Pi_9=223{,}092{,}870$. Six
candidate functions were tested and five died, among them the control $\delta=\log n$, which fails
first at $n=30$ where $\sigma=1.0333$; a battery that did not kill it would not be testing anything
\textup{(\texttt{code/invent\_lyapunov.rs})}.

That interval is an artefact of the range swept. No $\lambda$ works.

\begin{theorem}[The criterion fails, for every $\lambda$]\label{thm:lyapfalse}
There is no real $\lambda$ for which $(\ast)$ holds at every squarefree $n$ with $n'$ squarefree.
Two witnesses suffice. Let
\[ n_0=\!\!\!\prod_{\substack{7\le p\le373\\ p\notin\{307,\,317,\,359\}}}\!\!\!p, \]
a squarefree integer with $68$ prime factors and $142$ decimal digits. Then $n_0'=2\cdot3\cdot5\cdot C$
with $C$ a prime of $141$ digits, so $n_0'$ is squarefree and $\sigma(n_0')=\tfrac{31}{30}$ exactly,
while $\sigma(n_0)=1.0059067\ldots$. Since $\sigma(n_0)>1$ and $\sigma(n_0')>\sigma(n_0)$,
\[ \log\sigma(n_0)\;>\;0\;>\;\lambda\bigl(\sigma(n_0)-\sigma(n_0')\bigr)\qquad(\lambda>0), \]
so $(\ast)$ fails at $n_0$ for every positive $\lambda$. At $n=30$, with $30'=31$ prime,
$\sigma(30)=\tfrac{31}{30}>1$ and $\sigma(30)-\sigma(31)=\tfrac{931}{930}>0$, so $\log\sigma(30)>0\ge
\lambda(\sigma(30)-\sigma(31))$ for every $\lambda\le0$. Together the two exclude every real
$\lambda$.
\end{theorem}
\begin{proof}
Both statements are finite verifications in exact rational arithmetic. For $n_0$: $\sigma(n_0)$ and
$\sigma(n_0')$ are computed as rationals, $\gcd(n_0,n_0')=1$ and $\sigma(n_0)=n_0'/n_0$ confirm
Theorem~\ref{K-thm:structure}, the factorisation $n_0'=2\cdot3\cdot5\cdot C$ multiplies back to $n_0'$
and has all exponents $1$, and $C$ carries an Atkin--Morain certificate, generated and validated
\textup{(\texttt{data/certs/lyap\_refute\_cofactor\_ecpp.txt})}. The comparisons $\sigma(n_0)>1$ and
$\sigma(n_0')>\sigma(n_0)$ are exact, and $\sigma(n_0')=\sigma(n_0)$ is impossible in any case by
Lemma~\ref{C-lem:sigmainj}. The case $n=30$ is arithmetic
\textup{(\texttt{code/lyap\_refute\_verify.gp})}.
\end{proof}

The argument is formalised. \texttt{Erdos307.criterion\_fails\_of\_pos} and
\texttt{criterion\_fails\_of\_nonpos} give the two sign clashes,
\texttt{no\_lambda\_works} combines them, and \texttt{lyapunov\_criterion\_false} instantiates the
combination at $\sigma(n_0)$, carried as the exact quotient $n_0'/n_0$, and at the pair
$(\tfrac{31}{30},\tfrac1{31})$ of $n=30$; the two numeric comparisons $1<\sigma(n_0)$ and
$\sigma(n_0)<\tfrac{31}{30}$ are checked on the $142$-digit numerals themselves. All six
declarations carry $0$ \texttt{sorry} and only \texttt{propext}, \texttt{Classical.choice},
\texttt{Quot.sound}, with no \texttt{native\_decide}
\textup{(\texttt{lean/Erdos307/LyapFalse.lean})}. What Lean does not see is the factorisation
$n_0'=2\cdot3\cdot5\cdot C$ and the primality of $C$, which is the Atkin--Morain certificate; the
status is therefore \textup{VERIFIED} modulo that certificate, in the same sense that
Theorem~\ref{thm:halflyap} is \textup{VERIFIED} modulo Theorem~\ref{K-thm:structure}.

The witness is built, not found. For squarefree $n$ and a prime $w\nmid n$ one has
$n'\equiv n\sum_{p\mid n}p^{-1}\pmod w$, hence
\[ w\mid n'\quad\Longleftrightarrow\quad \sum_{p\mid n}p^{-1}\equiv0 \pmod w. \tag{F} \]
Imposing (F) for $w\in\{2,3,5\}$ forces $\sigma(n')\ge\tfrac{31}{30}$, and building $n$ from primes
$\ge7$ with $\sigma(n)$ just above $1$ gives $\sigma(n')>\sigma(n)$. Reaching mass $1$ without
$2,3,5$ needs primes only to about $360$, which is what keeps $n_0'$ small enough to factor
completely; the squarefreeness of $n_0'$ is therefore verified rather than assumed. Three of the
larger pool primes are dropped to satisfy (F), at negligible cost in mass
\textup{(\texttt{code/lyap\_refute.gp})}.

The sweep could not have seen this. Every $n\le4\times10^8$ with $\sigma(n)>1$ is even, so the
battery measured only $n$ containing $2$, whereas the failure needs $n$ to omit $2$, $3$ and $5$ so
that $n'$ may absorb them. By Theorem~\ref{thm:rhobarrier} below, any failure must have
$n\,n'\ge\Pi_{59}$, and $n_0n_0'$ has $284$ digits: the region was out of reach of any sweep, and
only a construction enters it.

One route to the criterion had looked provable rather than measured, and it fails for the same
reason. Write $\rho=\sup\{\sigma(n')/\sigma(n):\sigma(n)>1\}$, the supremum being over the same
class of $n$. If $\rho\le c<1$ then $\sigma(n)-\sigma(n')\ge(1-c)\sigma(n)$, and since
$\log x/x\le1/e$,
\[ L:=\sup_{\sigma(n)>1}\frac{\log\sigma(n)}{\sigma(n)-\sigma(n')}\;\le\;\frac1{e(1-c)} . \]
Measured over $n\le4\times10^8$, $\rho=0.519746$, which would give $L\le0.766$ against a measured
upper demand of $1.274$; a proof of $\rho<1-1/(eU)$, for $U$ a lower bound on the upper demand,
would then have given $(\ast)$. The measurement is not the truth. Already
\[ a_0=\!\!\prod_{5\le p\le163}\!\!p,\qquad
   a_0'=2\cdot3\cdot173\cdot227\cdot1973\cdot82890547\cdot68126695363\cdot P_{36}, \]
with $P_{36}$ prime, is squarefree with squarefree derivative, $\sigma(a_0)=1.0723027\ldots>1$ and
$\sigma(a_0')=0.8440258\ldots$, so $\rho\ge0.7871152\ldots$; the bound the route would yield is
$1/(e(1-\rho))\ge1.7280$, already above the measured upper demand, so it constrains nothing. The
witness $n_0$ of Theorem~\ref{thm:lyapfalse} gives $\rho\ge1.02726$, so $\rho<1$ is false outright
\textup{(\texttt{code/rho\_verify.gp})}.

The mechanism advanced for $\rho<1$ was that $\sigma(n)>1$ forces $n$ to absorb the small primes
while $n'$ stays coprime to $n$. It does not: the reciprocals of the primes diverge, so $\sigma(n)>1$
is reachable from primes $\ge7$ alone, and then $n'$ is free to take $2$, $3$ and $5$. Both $a_0$ and
$n_0$ are of that shape. The earlier probe of this regime evaluated $\sigma(n')/\sigma(n)$ at the
least odd $n$ with $\sigma(n)>1$, obtaining $0.0312$, and read the regime off that single point;
the quantity to be bounded is a supremum, and the least member of a family does not report it. This
is the error recorded in Remark~\ref{C-rem:sectorsquared}, in a second instance.

What survives is the location of the failure, and it is exactly where the barrier puts it.

\begin{theorem}[$\rho<1$ below the barrier]\label{thm:rhobarrier}
Let $n$ be squarefree with $n'$ squarefree and $\sigma(n)>1$. If $\sigma(n')\ge\sigma(n)$ then
$n\,n'\ge\Pi_{59}>8.77\times10^{112}$. Consequently $\rho<1$ for every $n$ with
$n\,n'<\Pi_{59}$, and since $\sigma(n)>1$ gives $n'>n$, for every $n<\sqrt{\Pi_{59}}$, that is for
every $n<10^{56.5}$.
\end{theorem}
\begin{proof}
$\sigma(n')\ge\sigma(n)>1$ gives $\sigma(n)+\sigma(n')>2$. By Theorem~\ref{K-thm:structure} $n$ and
$n'$ are coprime, and both are squarefree, so $\sigma(n\,n')=\sigma(n)+\sigma(n')>2$. A squarefree
integer of mass at least $2$ has at least $59$ prime factors and is therefore at least $\Pi_{59}$,
since $\sigma(\Pi_{58})=1.998740<2\le2.002350=\sigma(\Pi_{59})$; this is the mass estimate
underlying Theorem~\ref{K-thm:barrier}. Finally $\sigma(n)>1$ gives $n'=n\sigma(n)>n$, so
$n^2<n\,n'$.
\end{proof}

The counting step is formalised. \texttt{Erdos307.card\_ge\_59\_of\_recipSum\_ge\_two} isolates the
mass estimate, that a set of primes of reciprocal sum at least $2$ has at least $59$ elements, and
\texttt{Erdos307.card\_union\_ge\_59\_of\_masses} applies it to two disjoint sets each of mass
exceeding $1$; both carry $0$ \texttt{sorry}
\textup{(\texttt{lean/Erdos307/RhoBarrier.lean})}. Their axiom lists are those of
Theorem~\ref{K-thm:barrier} itself, the three standard axioms together with the
\texttt{native\_decide} evaluation of the first $58$ primes. The same core proves
\texttt{card\_ge\_59} of the barrier, which assumes $\sigma(P)\sigma(Q)=1$ instead; the two
hypotheses differ and the shared step has been extracted rather than duplicated.

Theorem~\ref{thm:rhobarrier} therefore locates the failure before exhibiting it: $\rho\ge1$ forces
$n\,n'\ge\Pi_{59}$, so no sweep can meet a counterexample, and any counterexample must have $n$ of
at least $57$ digits. The witness $n_0$ has $142$, and $n_0n_0'$ has $284$. The regime in which
$\rho\ge1$ is possible is the regime in which a two-cycle is possible, since both need two coprime
squarefree integers of total mass at least $2$; what Theorem~\ref{thm:lyapfalse} shows is that the
first of those is genuinely non-empty, which says nothing about the second.

\begin{corollary}[The half-mass hypothesis cannot be relaxed]\label{cor:halfsharp}
For every $\lambda>0$ and every $c>2.03925$, the implication
``$\sigma(n\,n')<c$ implies $\delta_\lambda(n')<\delta_\lambda(n)$'' is false. Indeed
$\sigma(n_0\,n_0')=2.039240\ldots$ while
$\delta_\lambda(n_0')-\delta_\lambda(n_0)=\log\sigma(n_0)+\lambda(\sigma(n_0')-\sigma(n_0))>0$ for
every $\lambda>0$.
\end{corollary}

Two is also the exact threshold for the argument of Theorem~\ref{thm:halflyap}, in the following
sense, which concerns the inequality alone and not the arithmetic of $\sigma$.

\begin{proposition}[The mass-sum relaxation stops at $2$]\label{prop:masssumthreshold}
Let $c>0$. There is $f:(0,\infty)\to\mathbb{R}$ with $f(x)-f(y)>\log x$ for all $x,y>0$ satisfying
$x+y<c$ if and only if $c\le2$. For $c=2$ the linear solutions are exactly $f(x)=\tfrac12x+\text{const}$.
\end{proposition}
\begin{proof}
The region $x+y<c$ is symmetric, so a solution gives both $f(x)-f(y)>\log x$ and $f(y)-f(x)>\log y$;
adding, $0>\log(xy)$. Taking $x=y\uparrow c/2$ forces $c^2/4\le1$, that is $c\le2$. For $c\le2$ take
$f(x)=\tfrac12x$: if $x+y<2$ then $\tfrac12(x-y)>x-1\ge\log x$. If $f=\lambda x$ then letting
$y\uparrow2-x$ gives $2\lambda(x-1)\ge\log x$ for all $x\in(0,2)$, and dividing by $x-1$ and letting
$x\to1$ from each side gives $2\lambda\ge1$ and $2\lambda\le1$.
\end{proof}

Both halves are formalised: \texttt{Erdos307.half\_works} for $c\le2$ and
\texttt{no\_f\_above\_two} for $c>2$, the latter by the diagonal point $x=y$ with $1<x<c/2$, where
the right side vanishes and the left is positive; \texttt{mass\_sum\_threshold} states the two
together. The uniqueness of the constant is formalised as well: \texttt{lambda\_eq\_half} shows that a linear solution at $c=2$ forces $\lambda=\tfrac12$, by the local-minimum route rather than a two-sided limit, since $g(x)=2\lambda(x-1)-\log x$ is nonnegative on $(0,2)$ and vanishes at $1$, so its derivative $2\lambda-1$ vanishes there \textup{(\texttt{lean/Erdos307/NoInvariant.lean})}.

Any strengthening must therefore use an input about $n$ beyond the pair
$(\sigma(n),\sigma(n'))$, and by the next proposition there is no room in the choice of function
either.

\begin{proposition}[The ansatz is universal, so the search was the problem]\label{prop:lyapuniversal}
Let $S$ be the set of squarefree $n\ge2$ with $n'$ squarefree. The following are equivalent:
\begin{enumerate}[label=\textup{(\arabic*)},leftmargin=2.2em]
\item $n\mapsto n'$ has no periodic orbit inside $S$;
\item there is $\delta:S\to\mathbb{R}$ with $\delta(n')<\delta(n)$ whenever $n,n'\in S$;
\item there is $f$ on $\sigma(S)$ with $f(\sigma(n))-f(\sigma(n'))>\log\sigma(n)$ for all such $n$.
\end{enumerate}
\end{proposition}
\begin{proof}
(2) implies (1), since a periodic orbit would give $\delta(n)<\delta(n)$. For (1) implies (2), the
transitive closure of $n\mapsto n'$ is irreflexive by (1) and transitive by construction, hence a
strict partial order on the countable set $S$; it has a linear extension, and a countable linear
order embeds order-preservingly in $\mathbb{Q}$, which supplies $\delta$. For (2) and (3), $\sigma$
is injective on squarefree integers by Lemma~\ref{C-lem:sigmainj}, so every $\delta$ on $S$ is
$\log n+f(\sigma(n))$ for exactly one $f$, namely $f(\sigma(n))=\delta(n)-\log n$; and $n'=n\sigma(n)$
turns $\delta(n')<\delta(n)$ into the inequality of (3).
\end{proof}

The collapse of the ansatz is formalised. \texttt{Erdos307.exists\_f\_of\_injective} proves that
injectivity of $\sigma$ forces \emph{every} $\delta$ into the shape $\log n+f(\sigma(n))$,
\texttt{increment\_iff} converts $\delta(n')<\delta(n)$ into $(\ast)$ verbatim, and
\texttt{no\_two\_cycle\_of\_decreasing} is the implication of
Proposition~\ref{prop:lyapworks}; all with $0$ \texttt{sorry} and the three standard axioms
\textup{(\texttt{lean/Erdos307/NoInvariant.lean})}.

So Definition~\ref{def:lyap} restricts nothing. The tournament that produced $\delta_\lambda$ was
searching the space of all functions on $S$, and by Proposition~\ref{prop:lyapuniversal} the
existence of any member of that space is equivalent to the absence of cycles, which is strictly
stronger than a negative answer to \#307. The criterion was never a weakening of the problem, and
Theorem~\ref{thm:lyapfalse} settles the one form of it that had been left open. What remains true,
and is the residue of the whole construction, is Theorem~\ref{thm:halflyap} together with
Corollary~\ref{cor:lyapbarrier}: an elementary Lyapunov proof of the barrier, valid exactly up to
mass $2$ and, by Corollary~\ref{cor:halfsharp}, not one step further.

\emph{Three independent blocks on $\mathrm{A}10$.} The existence route is therefore obstructed at three separate levels, and no two of them are the
same obstruction. Proposition~\ref{prop:nofreeslot} removes the freedom: the construction carries
one parameter and cannot be relaxed into a denser family. Theorem~\ref{thm:ff} removes the
method: no argument that survives the passage to $K[t]$ can prove existence, because there the
statement is false. Proposition~\ref{prop:multiplicityfail} removes the last refuge, the hope that a
large family of sequences is easier than one: the union stays logarithmically thin however many
sequences it has. A fourth block was claimed here and does not stand: it held that the $34$ immune
families could not be examined, and Proposition~\ref{A-prop:immunedecide} examines them and finds them
empty (Remark~\ref{rem:untestable}). Removing it strengthens nothing and is recorded because the
count of independent obstructions is exactly the kind of claim that should not drift.
What is left is a lower bound for primes in a sparse exponential sequence, which exists for no
non-degenerate linear recurrence at all. We state plainly that the present work does not decide the
existence direction and does not contain a route to deciding it.



\begin{thebibliography}{9}
\bibitem{Barbeau1961} E.~J.~Barbeau, \emph{Remarks on an arithmetic derivative}, Canad.\ Math.\
Bull.\ \textbf{4} (1961), 117--122.
\bibitem{Barbeau1976} E.~J.~Barbeau, \emph{Computer challenge corner: Problem 477}, J.\ Recreational
Math.\ \textbf{9} (1976/77), 30.
\bibitem{ErdosGraham1980} P.~Erd\H{o}s and R.~L.~Graham, \emph{Old and New Problems and Results in
Combinatorial Number Theory}, Monogr.\ Enseign.\ Math.\ \textbf{28}, Geneva, 1980, p.~38.
\bibitem{GrahamEgyptian} R.~L.~Graham, \emph{Paul Erd\H{o}s and Egyptian fractions}, in \emph{Erd\H{o}s
Centennial}, Bolyai Soc.\ Math.\ Stud.\ \textbf{25}, Springer, 2013, pp.~289--309.
\bibitem{Watanabe2020} T.~Watanabe, \emph{New examples of the representation of $1$ by the sum of
reciprocals of semiprime numbers}, arXiv:2009.03275 (2020).
\bibitem{Bado2026b} I.-O.~Bado, \emph{Primary-pseudoperfect sectors in two-cycles of the
arithmetic derivative}, preprint (2026), communicated to the author.
\bibitem{BloomShifted} T.~F.~Bloom, \emph{Unit fractions with shifted prime denominators}, Proc.\
Roy.\ Soc.\ Edinburgh Sect.\ A (2024); arXiv:2305.02689.
\bibitem{UfnarovskiAhlander2003} V.~Ufnarovski and B.~\AA hlander, \emph{How to differentiate a
number}, J.\ Integer Seq.\ \textbf{6} (2003), Article 03.3.4.
\bibitem{Kovic2012} J.~Kovi\v{c}, \emph{The arithmetic derivative and antiderivative}, J.\ Integer
Seq.\ \textbf{15} (2012), Article 12.3.8.
\bibitem{MO320838} MathOverflow, \emph{Can the product of two sums of prime reciprocals be $1$?},
question 320838 (comments by R.~Israel and J.~Rosen, answer by ``Woett''), 2019.
\url{https://mathoverflow.net/questions/320838}.
\bibitem{ErdosProblems307} T.~F.~Bloom, \emph{Erd\H{o}s problem \#307},
\url{https://www.erdosproblems.com/307}.
\bibitem{OEIS_A003415} OEIS Foundation, \emph{Sequence A003415 (arithmetic derivative)},
\url{https://oeis.org/A003415}.
\bibitem{Pollack2018} P.~Pollack, \emph{On amicable tuples}, Illinois J. Math. \textbf{62} (2018),
  no.~1--4, 225--240; arXiv:1711.06847. (Explicit bounds for relatively prime amicable pairs,
  refining the finiteness of Artjuhov and Borho.)

\bibitem{teRiele1984} H.~J.~J.~te~Riele, \emph{Computation of all the amicable pairs below
$10^{10}$}, Math.\ Comp.\ \textbf{47} (1986), 361--368.
\bibitem{BHP2001} R.~C.~Baker, G.~Harman, and J.~Pintz, \emph{The difference between consecutive
primes, II}, Proc.\ London Math.\ Soc.\ (3) \textbf{83} (2001), 532--562.
\bibitem{Iwaniec1978} H.~Iwaniec, \emph{Almost-primes represented by quadratic polynomials},
Invent.\ Math.\ \textbf{47} (1978), 171--188.
\bibitem{LemkeOliver2012} R.~J.~Lemke Oliver, \emph{Almost-primes represented by quadratic
polynomials}, Acta Arith.\ \textbf{151} (2012), 241--261.
\bibitem{HardyWright} G.~H.~Hardy and E.~M.~Wright, \emph{An Introduction to the Theory of Numbers},
6th ed., Oxford Univ.\ Press, 2008.
\bibitem{Tenenbaum} G.~Tenenbaum, \emph{Introduction to Analytic and Probabilistic Number Theory},
3rd ed., Grad.\ Stud.\ Math.\ \textbf{163}, Amer.\ Math.\ Soc., 2015 (Erd\H{o}s--Wintner theorem).
\bibitem{Stothers} W.~W.~Stothers, \emph{Polynomial identities and Hauptmoduln}, Quart.\ J.\ Math.\
Oxford (2) \textbf{32} (1981), 349--370; R.~C.~Mason, \emph{Diophantine Equations over Function
Fields}, Cambridge Univ.\ Press, 1984.
\bibitem{Seva323194} Seva (\texttt{mathoverflow.net} user), \emph{A recursion with a number--theoretic
function}, MathOverflow question 323194 (2019). \url{https://mathoverflow.net/questions/323194}.
\bibitem{FewPrimes2025} \emph{Egyptian fractions for few primes}, arXiv:2507.03727 (July 2025).
\bibitem{Steinerberger2025} S.~Steinerberger, \emph{On an iterated arithmetic function problem of
Erd\H{o}s and Graham}, arXiv:2504.08023 (April 2025).
\bibitem{StretchedEG2026} \emph{A stretched--exponential bound for an Erd\H{o}s--Graham
unit--fraction problem}, arXiv:2607.04157 (July 2026).
\bibitem{PortFillings2026} \emph{Port fillings for primary pseudoperfect numbers}, arXiv:2605.21518
(May 2026).
\bibitem{Semiprime2026} \emph{Every natural number is a sum of distinct semiprime unit fractions},
arXiv:2606.15159 (June 2026).
\bibitem{ButskeJajeMayernik} W.~Butske, L.~M.~Jaje and D.~R.~Mayernik, \emph{On the equation
$\sum_{p\mid N}1/p+1/N=1$, pseudoperfect numbers, and perfectly weighted graphs}, Math.\ Comp.\
\textbf{69} (2000), 407--420.
\bibitem{Bado2026} I.~O.~Bado, \emph{Structural constraints for an Erd\H{o}s unit--fraction problem
over primes}, preprint, Aix--Marseille University (May 2026).
\url{https://doi.org/10.13140/RG.2.2.32245.54245}.
\bibitem{Giuga1950} G.~Giuga, \emph{Su una presumibile propriet\`a caratteristica dei numeri primi},
Ist.\ Lombardo Sci.\ Lett.\ Rend.\ Cl.\ Sci.\ Mat.\ Nat.\ (3) \textbf{14} (1950), 511--528.
\bibitem{BBBG1996} D.~Borwein, J.~M.~Borwein, P.~B.~Borwein, and R.~Girgensohn, \emph{Giuga's conjecture
on primality}, Amer.\ Math.\ Monthly \textbf{103} (1996), no.~1, 40--50.
\bibitem{BorweinWong1997} J.~M.~Borwein and E.~Wong, \emph{A survey of results relating to Giuga's
conjecture on primality}, CRM Proc.\ Lecture Notes \textbf{11} (1997), 13--27.
\bibitem{BJM2000} W.~Butske, L.~M.~Jaje, and D.~R.~Mayernik, \emph{On the equation
$\sum_{p\mid N}1/p+1/N=1$, pseudoperfect numbers, and perfectly weighted graphs}, Math.\ Comp.\
\textbf{69} (2000), no.~229, 407--420.
\bibitem{GrauOllerMarcen2012} J.~M.~Grau and A.~M.~Oller-Marc\'en, \emph{Giuga numbers and the
arithmetic derivative}, J.\ Integer Seq.\ \textbf{15} (2012), Article 12.4.1; arXiv:1103.2298.
\bibitem{SondowMacMillan2017} J.~Sondow and K.~MacMillan, \emph{Primary pseudoperfect numbers,
arithmetic progressions, and the Erd\H{o}s--Moser equation}, Amer.\ Math.\ Monthly \textbf{124} (2017),
no.~3, 232--240; arXiv:1812.06566.
\bibitem{GrauOllerMarcenSadornil2021} J.~M.~Grau, A.~M.~Oller-Marc\'en, and D.~Sadornil,
\emph{On $\mu$-Sondow numbers}, arXiv:2111.14211 (2021); Acta Math.\ Hungar.\ (2022).
\bibitem{Cox} D.~A.~Cox, \emph{Primes of the Form $x^2+ny^2$: Fermat, Class Field Theory, and Complex
Multiplication}, 2nd ed., Wiley, 2013.
\bibitem{OEIS_A054377} OEIS Foundation, \emph{Sequence A054377 (primary pseudoperfect numbers)},
\url{https://oeis.org/A054377}.
\bibitem{OEIS_A007850} OEIS Foundation, \emph{Sequence A007850 (Giuga numbers)},
\url{https://oeis.org/A007850}.
\bibitem{OEIS_A396915} OEIS Foundation, \emph{Sequence A396915 (composite squarefree $k$ with
$k'+2k$ a perfect square)}, \url{https://oeis.org/A396915}.
\bibitem{Evertse1984} J.-H.~Evertse, \emph{On sums of $S$-units and linear recurrences}, Compositio
Math.\ \textbf{53} (1984), 225--244.
\bibitem{FI1998} J.~Friedlander and H.~Iwaniec, \emph{The polynomial $X^2+Y^4$ captures its
primes}, Ann. of Math. \textbf{148} (1998), 945--1040.
\bibitem{HB2001} D.~R.~Heath-Brown, \emph{Primes represented by $x^3+2y^3$}, Acta Math.
\textbf{186} (2001), 1--84.
\bibitem{Stewart2013} C.~L.~Stewart, \emph{On divisors of Lucas and Lehmer numbers}, Acta Math.
\textbf{211} (2013), 291--314.
\bibitem{BCZ2003} Y.~Bugeaud, P.~Corvaja, and U.~Zannier, \emph{An upper bound for the G.C.D. of
$a^n-1$ and $b^n-1$}, Math. Z. \textbf{243} (2003), 79--84.
\bibitem{Merikoski2019} J.~Merikoski, \emph{On the largest prime factor of $n^2+1$}, J. Eur. Math.
Soc. \textbf{25} (2023), 1253--1284.
\bibitem{BGS2010} J.~Bourgain, A.~Gamburd, and P.~Sarnak, \emph{Affine linear sieve, expanders,
and sundry applications}, Invent. Math. \textbf{179} (2010), 559--644.
\bibitem{ESS2002} J.-H.~Evertse, H.~P.~Schlickewei, and W.~M.~Schmidt, \emph{Linear equations in
variables which lie in a multiplicative group}, Ann.\ of Math.\ (2) \textbf{155} (2002), 807--836.
\bibitem{Erdos1952} P.~Erd\H{o}s, \emph{On the sum $\sum_{k=1}^x d(f(k))$}, J.\ London Math.\ Soc.\
\textbf{27} (1952), 7--15.
\bibitem{BHV2001} Y.~Bilu, G.~Hanrot, and P.~M.~Voutier, \emph{Existence of primitive divisors of
Lucas and Lehmer numbers} (with an appendix by M.~Mignotte), J.\ Reine Angew.\ Math.\ \textbf{539}
(2001), 75--122.
\end{thebibliography}
\end{document}
